\documentclass[11pt]{article}

\usepackage[letterpaper,margin=0.82in,headheight=15pt]{geometry}
\usepackage[T1]{fontenc}
\usepackage{lmodern}
\usepackage{microtype}
\usepackage{amsmath,amssymb,amsthm,bm,mathtools}
\usepackage{booktabs,longtable,tabularx,array,multirow}
\usepackage{enumitem}
\usepackage{xcolor}
\usepackage{fancyhdr}
\usepackage{graphicx}
\usepackage{tikz-cd}
\usepackage{hyperref}
\usepackage[nameinlink,noabbrev]{cleveref}

\definecolor{navy}{HTML}{17324D}
\definecolor{teal}{HTML}{147D92}
\definecolor{orange}{HTML}{D47A27}
\definecolor{softblue}{HTML}{EEF4F7}
\definecolor{softgray}{HTML}{F3F5F6}
\definecolor{darkgray}{HTML}{46545E}
\definecolor{softorange}{HTML}{FFF4E8}
\definecolor{softgreen}{HTML}{EDF7F0}
\definecolor{softred}{HTML}{FCEEEE}

\hypersetup{
  colorlinks=true,
  linkcolor=navy,
  citecolor=teal,
  urlcolor=teal,
  pdftitle={Volume VII: Concrete Microscopic Nucleon Hamiltonian and Renormalization Program},
  pdfauthor={DeuteronWigner project}
}

\pagestyle{fancy}
\fancyhf{}
\fancyhead[L]{\small\color{darkgray}DeuteronWigner formalism - Volume VII}
\fancyhead[R]{\small\color{darkgray}Microscopic Hamiltonian and renormalization}
\fancyfoot[L]{\small\color{darkgray}30 July 2026}
\fancyfoot[R]{\small\color{darkgray}\thepage}
\renewcommand{\headrulewidth}{0.4pt}

\setlist[itemize]{leftmargin=1.5em,itemsep=2pt,topsep=3pt}
\setlist[enumerate]{leftmargin=1.7em,itemsep=2pt,topsep=3pt}
\setlength{\parindent}{0pt}
\setlength{\parskip}{5pt}
\setlength{\emergencystretch}{2em}
\renewcommand{\arraystretch}{1.16}
\allowdisplaybreaks

\newtheorem{definition}{Definition}[section]
\newtheorem{axiom}[definition]{Axiom}
\newtheorem{principle}[definition]{Design principle}
\newtheorem{requirement}[definition]{Requirement}
\newtheorem{proposition}[definition]{Proposition}
\newtheorem{theorem}[definition]{Theorem}
\newtheorem{lemma}[definition]{Lemma}
\newtheorem{corollary}[definition]{Corollary}
\newtheorem{remark}[definition]{Remark}
\newtheorem{decision}[definition]{Model decision}

\crefname{definition}{definition}{definitions}
\crefname{axiom}{axiom}{axioms}
\crefname{principle}{principle}{principles}
\crefname{requirement}{requirement}{requirements}
\crefname{proposition}{proposition}{propositions}
\crefname{theorem}{theorem}{theorems}
\crefname{lemma}{lemma}{lemmas}
\crefname{corollary}{corollary}{corollaries}
\crefname{remark}{remark}{remarks}
\crefname{decision}{model decision}{model decisions}

\newcommand{\kT}{\bm{k}_{T}}
\newcommand{\pT}{\bm{p}_{T}}
\newcommand{\qT}{\bm{q}_{T}}
\newcommand{\DT}{\bm{\Delta}_{T}}
\newcommand{\bD}{\bm{b}_{\Delta}}
\newcommand{\bM}{\bm{b}_{\mathrm{TMD}}}
\newcommand{\dd}{\mathrm{d}}
\newcommand{\Tr}{\operatorname{Tr}}
\newcommand{\tr}{\operatorname{tr}}
\newcommand{\diag}{\operatorname{diag}}
\newcommand{\rank}{\operatorname{rank}}
\newcommand{\Span}{\operatorname{span}}
\newcommand{\Ker}{\operatorname{Ker}}
\newcommand{\Ran}{\operatorname{Ran}}
\newcommand{\Cov}{\operatorname{Cov}}
\newcommand{\Var}{\operatorname{Var}}
\newcommand{\Id}{\operatorname{Id}}
\newcommand{\Hil}{\mathcal H}
\newcommand{\LF}{\ensuremath{\mathrm{LF}}}
\newcommand{\BLFQ}{\ensuremath{\mathrm{BLFQ}}}
\newcommand{\TMD}{\ensuremath{\mathrm{TMD}}}
\newcommand{\GTMD}{\ensuremath{\mathrm{GTMD}}}
\newcommand{\GPD}{\ensuremath{\mathrm{GPD}}}
\newcommand{\QCD}{\ensuremath{\mathrm{QCD}}}
\newcommand{\EFT}{\ensuremath{\mathrm{EFT}}}
\newcommand{\RGPEP}{\ensuremath{\mathrm{RGPEP}}}
\newcommand{\SRG}{\ensuremath{\mathrm{SRG}}}
\newcommand{\Amp}{\mathbf{Amp}}
\newcommand{\Dens}{\mathbf{Dens}}
\newcommand{\Match}{\mathbf{Match}}
\newcommand{\Red}{\mathbf{Red}}
\newcommand{\Proc}{\mathbf{Proc}}
\newcommand{\Infer}{\mathbf{Infer}}
\newcommand{\cO}{\mathcal O}
\newcommand{\cM}{\mathcal M}
\newcommand{\cR}{\mathcal R}
\newcommand{\cS}{\mathcal S}
\newcommand{\cV}{\mathcal V}
\newcommand{\cW}{\mathcal W}
\newcommand{\cE}{\mathcal E}
\newcommand{\cG}{\mathcal G}
\newcommand{\cK}{\mathcal K}
\newcommand{\cQ}{\mathcal Q}
\newcommand{\cP}{\mathcal P}
\newcommand{\cC}{\mathcal C}
\newcommand{\cZ}{\mathcal Z}
\newcommand{\cD}{\mathcal D}
\newcommand{\cN}{\mathcal N}
\newcommand{\cU}{\mathcal U}
\newcommand{\cF}{\mathcal F}
\newcommand{\cA}{\mathcal A}
\newcommand{\cB}{\mathcal B}
\newcommand{\cL}{\mathcal L}
\newcommand{\cJ}{\mathcal J}
\newcommand{\cI}{\mathcal I}
\newcommand{\cH}{\mathcal H}
\newcommand{\cT}{\mathcal T}
\newcommand{\cX}{\mathcal X}
\newcommand{\cY}{\mathcal Y}
\newcommand{\eps}{\varepsilon}
\newcommand{\bbone}{\bm 1}
\newcommand{\abs}[1]{\left|#1\right|}
\newcommand{\norm}[1]{\left\lVert#1\right\rVert}
\newcommand{\inner}[2]{\left\langle #1\middle|#2\right\rangle}
\newcommand{\avg}[1]{\left\langle #1\right\rangle}
\newcommand{\trans}{^{\mathsf T}}
\newcommand{\phys}{\mathrm{phys}}
\newcommand{\eff}{\mathrm{eff}}
\newcommand{\bare}{\mathrm{bare}}
\newcommand{\ind}{\mathrm{ind}}
\newcommand{\ct}{\mathrm{ct}}
\newcommand{\val}{\mathrm{val}}
\newcommand{\sea}{\mathrm{sea}}
\newcommand{\gl}{\mathrm{g}}
\newcommand{\reg}{\mathrm{reg}}
\newcommand{\trunc}{\mathrm{trunc}}
\newcommand{\conf}{\mathrm{conf}}
\newcommand{\chiop}{\chi}
\newcommand{\inst}{\mathrm{inst}}
\newcommand{\can}{\mathrm{can}}
\newcommand{\ren}{\mathrm{ren}}
\newcommand{\num}{\mathrm{num}}
\newcommand{\calib}{\mathrm{cal}}
\newcommand{\hold}{\mathrm{hold}}
\newcommand{\pilot}{\mathrm{pilot}}
\newcommand{\UV}{\mathrm{UV}}
\newcommand{\IR}{\mathrm{IR}}
\newcommand{\COM}{\mathrm{CM}}
\newcommand{\Poincare}{Poincar\'e}
\newcommand{\FSDR}{\ensuremath{\mathrm{FSDR}}}
\newcommand{\EMT}{\ensuremath{\mathrm{EMT}}}
\newcommand{\PCAC}{\ensuremath{\mathrm{PCAC}}}
\newcommand{\WTI}{\ensuremath{\mathrm{WTI}}}
\newcommand{\Lz}{L_z}
\newcommand{\Jz}{J^z}
\newcommand{\PauliLub}{W^\mu}
\newcommand{\StateBundle}{\texttt{MicroscopicStateBundle}}
\newcommand{\Ham}{\texttt{MicroscopicHamiltonian}}
\newcommand{\RenDatum}{\texttt{RenormalizationDatum}}
\newcommand{\TruncTower}{\texttt{HamiltonianTruncationTower}}

\newcolumntype{Y}{>{\raggedright\arraybackslash}X}
\newcolumntype{L}[1]{>{\raggedright\arraybackslash}p{#1}}

\title{\vspace{-1.0cm}
{\color{navy}\bfseries Volume VII: Concrete Microscopic Nucleon\\[2mm]
Hamiltonian and Renormalization Program}\\[4mm]
\large Selection of the primary light-front dynamical realization, sector-coupled QCD Hamiltonian,\\
counterterm flow, and the first state bundle admissible for common GTMD prediction}
\author{DeuteronWigner project}
\date{30 July 2026}

\begin{document}
\maketitle

\begin{center}
\begin{minipage}{0.94\textwidth}
\colorbox{softblue}{\parbox{0.96\textwidth}{
\textbf{Status and purpose.}
Volumes~I--VI specify the admissible mathematical architecture, regulated
state spaces, common nucleon overlap operators, Wilson-line phases, matched
spin-1 nuclear dynamics, regulator-to-QCD matching, evolution, process
factorization, and shared inference.  The present volume makes the first
irreversible dynamical choice.  It adopts a renormalized,
symmetry-adapted basis light-front Hamiltonian effective theory as the primary
microscopic route for the nucleon.  The model is defined by a resolution-indexed
Hamiltonian, explicit \(qqq\), \(qqqg\), and \(qqqq\bar q\) sectors,
sector-coupled counterterms and induced operators, exact color and permutation
constraints, a declared zero-mode remainder, and observable-level convergence
under enlargement of the basis, Fock space, and ultraviolet resolution.
Existing analytic C3--C6 states remain validation fixtures.  They do not become
microscopic predictions merely because their algebra closes.  The first
admissible output of this volume is a versioned proton/neutron state bundle
whose currents, operators, member identity, phases, and convergence manifest
are sufficiently complete to feed the common GTMD, Wilson, nuclear, matching,
and inference layers without introducing a new phenomenological function per
observable.
}}
\end{minipage}
\end{center}

\begin{abstract}
The DeuteronWigner program requires a microscopic state that generates quark,
antiquark, and gluon correlators from one renormalized Hamiltonian rather than
from separately normalized parton distributions.  This volume selects a
concrete realization: a symmetry-adapted basis light-front quantization
Hamiltonian effective theory with explicit \(qqq\), \(qqqg\), and
\(qqqq\bar q\) sectors, canonical light-front QCD vertices and instantaneous
interactions, resolution-dependent induced confinement and chiral operators,
and Fock-sector-dependent renormalization.  The transverse basis is a
center-of-mass-factorized two-dimensional harmonic oscillator basis; the
longitudinal basis is a positive-momentum discretization with an explicit
endpoint and zero-mode policy.  Exact color-singlet projection, full
identical-fermion antisymmetry, target \(J^z\), flavor, parton helicity, and
orbital labels are enforced before diagonalization.  The Hamiltonian,
currents, GTMD operators, and Wilson insertions are transformed together under
sector elimination and resolution change.  Renormalization conditions are
imposed along a directed truncation tower rather than at one convenient matrix
size, and convergence is claimed only for matched observables after
counterterm refitting.  The document defines the block Hamiltonian, regulator
set, counterterm ownership, chiral sea mechanism, confinement ablations,
Ward-identity and rotational-restoration conditions, eigensolver and state
tracking, calibration/withheld partitions, uncertainty components, software
interfaces, analytic benchmarks, negative tests, and final acceptance gates.
The resulting state bundle is designed to be the first central microscopic
parent from which the common GTMD and Wilson-line calculations can proceed.
\end{abstract}

\tableofcontents

\section{The dynamical decision}
\label{sec:decision}

\subsection{Why a new decision is required}

The existing phenomenological boundary has already demonstrated that a common
correlator contract can preserve flavor, helicity, transverse rank, Wilson-link
class, nuclear mechanism, and uncertainty provenance.  The analytic C3--C6
pilots have additionally demonstrated exact recoil maps, common regulated
reductions, explicit sea and gluon slots, Feshbach operator equivalence, and a
one-gluon absorptive cut.  Those successes are necessary but not sufficient.
Their wave functions are chosen analytic fixtures; their widths and sector
probabilities are not eigenstate predictions.  A microscopic calculation must
instead solve
\begin{equation}
 \cM_{\LF,r}^{2}(\theta_r)
 |\Psi_{N,\Lambda}^{(r)}\rangle
 =M_{N,r}^{2}|\Psi_{N,\Lambda}^{(r)}\rangle,
 \label{eq:master-eigenproblem}
\end{equation}
for a directed family of regulator and truncation labels \(r\), while
transforming every current and partonic operator consistently with the same
Hamiltonian.

\begin{principle}[State before distribution]
A named PDF, TMD, GPD, GTMD, form factor, or OAM quantity is not a Hamiltonian
parameter.  It is a reduction of a matrix element evaluated in a common
renormalized state.  No new coefficient may be introduced at the named-function
level unless it descends from a missing Hamiltonian, current, Wilson, matching,
or controlled-discrepancy operator.
\label{pr:state-before-distribution}
\end{principle}

\subsection{Primary route}

\begin{decision}[Renormalized symmetry-adapted BLFQ Hamiltonian EFT]
The primary nucleon realization is a \emph{renormalized,
symmetry-adapted basis light-front Hamiltonian effective theory} (RSA-BLFQ).
The basis and sparse eigensolver are BLFQ.  The operator content is a
resolution-dependent light-front QCD Hamiltonian supplemented by induced
confining, chiral, and truncation-restoration operators.  The coefficients of
those induced operators flow with the resolution and Fock content and are
fixed by shared renormalization conditions.  They are not permanent
phenomenological potentials whose parameters are copied unchanged between
truncations.
\label{dec:primary-route}
\end{decision}

The choice is motivated by five requirements that must hold simultaneously:
\begin{enumerate}
 \item explicit quark, antiquark, and gluon degrees of freedom in one state;
 \item sparse resolution of color, helicity, \(\Jz\), and \(\Lz\) blocks;
 \item a direct route to nonzero-transfer wave-function overlaps;
 \item explicit sector-changing QCD vertices needed by the Wilson and genuine
 quark--gluon sectors;
 \item a regulator and basis tower on which counterterm flow and convergence can
 be measured.
\end{enumerate}
This combination builds on BLFQ's sparse Hamiltonian program and its explicit
valence, dynamical-gluon, and light-sea extensions
\cite{Vary2010,Xu2021,Xu2023,Xu2025,Mondal2025}.

The word ``effective'' is indispensable.  At any finite Fock and basis
truncation, the retained Hamiltonian contains induced interactions that
represent excluded modes and sectors.  Predictivity comes from a controlled
flow and successful withheld predictions, not from calling a finite matrix the
canonical QCD Hamiltonian.

\subsection{Candidate-route comparison}

\begin{longtable}{L{0.19\textwidth}L{0.25\textwidth}L{0.25\textwidth}L{0.20\textwidth}}
\toprule
Route & Strength for this project & Principal obstruction & Assigned role\\
\midrule
\endhead
Valence-only confining BLFQ & Mature sparse basis, spin-resolved valence
amplitudes, useful current and GTMD benchmarks & No explicit antiquarks or
dynamical gluons; sea and Wilson phases would be external & Regression and
initial Hamiltonian benchmark only\\
Canonical-sector BLFQ without explicit confinement & Explicit
\(qqq\oplus qqqg\oplus qqqq\bar q\) QCD interactions and a direct common state
& At currently accessible finite resolutions, large effective masses,
incomplete antisymmetry, zero-mode omission, and limited convergence can absorb
missing confinement and higher sectors & Essential ablation and target of the
resolution flow; not accepted uncritically as the final finite Hamiltonian\\
RSA-BLFQ Hamiltonian EFT & Explicit partons plus controlled induced operators,
sector-dependent renormalization, sparse symmetry blocks, and a convergence
tower & Requires a disciplined counterterm basis, shared calibration, and
expensive multi-axis convergence & \textbf{Primary route}\\
Hadronic light-front Hamiltonian EFT & Natural chiral pion cloud, sea-flavor
asymmetry, and later nuclear matching & Hadronic degrees of freedom do not by
themselves generate the common quark/gluon GTMD parent & Matching and
long-distance chiral constraint on the partonic Hamiltonian\\
Dyson--Schwinger/Bethe--Salpeter amplitudes mapped to the light front &
Continuum covariance and strong chiral dynamics & Explicit Fock bookkeeping,
ordered Wilson phases, and sector-dependent operator matching are indirect &
Independent comparison and possible matching input\\
Light-front coupled cluster, RGPEP, or similarity flow & Potentially avoids hard
Fock cutoffs or generates systematic effective Hamiltonians & Full nucleon
quark/gluon implementation at the required spin and operator depth is not yet
the shortest path & Future solver or effective-Hamiltonian generator\\
Direct lattice-to-LF reconstruction & First-principles Euclidean matrix
elements and controlled continuum extrapolation & Real-time Wilson phases and
full LFWFs are not directly supplied & Calibration and withheld validation of
moments and matched distributions\\
\bottomrule
\end{longtable}

\begin{requirement}[No deeper patchwork]
Wave functions or interaction terms from different microscopic routes may not
be assembled into one central state unless an explicit matching map, overlap
subtraction, and shared renormalization datum are supplied.  Alternative routes
remain separate microscopic members or validation calculations.
\label{req:no-deeper-patchwork}
\end{requirement}

\section{Inherited contracts and scope}
\label{sec:scope}

\subsection{Inputs from Volumes I--VI}

Volume~VII inherits rather than redefines the following contracts:
\begin{enumerate}
 \item the light-front convention
 \(v^{\pm}=(v^0\pm v^3)/\sqrt2\) and
 \(\cM^2=2P^+P^- -\bm P_T^2\);
 \item the graded physical Hilbert spaces and distinct map classes
 \(\Amp,\Dens,\Match,\Red,\Proc,\Infer\);
 \item the zero-skewness momentum-fiber and recoil definitions;
 \item decorated quark, antiquark, and ordered-link gluon operator identities;
 \item common GTMD-to-TMD/GPD/PDF/current reductions;
 \item explicit Wilson-line pole, cut, color, and phase-budget identities;
 \item matched quark--gluon-to-nuclear operator separation;
 \item regulator-to-QCD matching and two-scale evolution as later maps;
 \item shared-parameter inference and frozen withheld-validation rules.
\end{enumerate}

The present volume specifies the dynamical object that must satisfy these
interfaces.  It does not change the accepted forward phenomenological parent
and does not promote any C3--C6 validation fixture to production.

\subsection{Declared initial physical scope}

The first microscopic state bundle targets the proton and neutron at a low
resolution \(\Lambda_h\), at zero skewness for initial GTMD tests, with light
flavors \(u,d\) and explicit transverse gluons.  The retained Fock space is
\begin{equation}
 \cF_{\mathrm{initial}}
 =\bigl\{
 |qqq\rangle,
 |qqqg\rangle,
 |qqqq\bar q\rangle_{q=u,d}
 \bigr\}.
 \label{eq:initial-fock-space}
\end{equation}
Strange pairs, two-gluon sectors, pair-plus-gluon sectors, and zero-mode
resolved sectors enter later truncation levels.  Their omission is represented
by a typed remainder and induced-operator basis, not by the assertion that they
are physically absent.

\begin{requirement}[Minimum orbital support]
Every retained sector must contain all basis blocks required for
\(\Lz=0,\pm1,\pm2\) interference after total-\(\Jz\) coupling.  A truncation
that removes a required block may be used as a controlled zero test, but not as
a complete prediction for Sivers, Boer--Mulders, worm-gear, pretzelosity,
linearly polarized gluon, or spin-1 tensor observables.
\label{req:minimum-oam}
\end{requirement}

\subsection{Explicit nonclaims}

This volume does not claim:
\begin{itemize}
 \item an exact solution of continuum QCD;
 \item that a two-dimensional harmonic-oscillator basis is physically
 privileged;
 \item that an induced confining operator survives unchanged in the continuum;
 \item that finite Fock probabilities are scheme-independent observables;
 \item that a low-resolution state is already a soft-subtracted TMD;
 \item that one fitted nucleon mass determines all counterterms;
 \item that a successful finite matrix spectrum establishes Poincar\'e
 covariance;
 \item that the first state bundle is ready for nuclear composition, evolution,
 or global inference before its corresponding gates pass.
\end{itemize}

\section{Resolution-indexed Hamiltonian}
\label{sec:hamiltonian}

\subsection{Master decomposition}

At truncation and regulator label
\begin{equation}
 r=\bigl(
 K,N_{\max},b,\lambda,\epsilon_x,\cF_r,
 \cB_r,\cR_r
 \bigr),
 \label{eq:regulator-label}
\end{equation}
the mass-squared operator is
\begin{align}
 \cM_{\LF,r}^{2}
 ={}&\cM_{0,r}^{2}
 +V_{qg,r}^{\can}
 +V_{3g,r}^{\can}
 +V_{4g,r}^{\can}
 +V_{\inst,r}^{\can}
 \nonumber\\
 &+V_{\conf,r}^{\ind}
 +V_{\chi,r}^{\ind}
 +V_{\rm spin,r}^{\ind}
 +\sum_{a}c_{a}^{(r)}\cO_{a}^{(r)}
 +\Delta\cM_{r}^{2,\trunc}.
 \label{eq:master-hamiltonian}
\end{align}
The first line contains canonical light-front QCD terms that act within the
retained sectors.  The second line contains induced low-resolution operators,
counterterms, and a quantified omitted-operator remainder.

\begin{definition}[Hamiltonian identity]
A Hamiltonian instance is identified by
\begin{equation}
 \mathfrak h_r=
 \left(
 \begin{gathered}
 \text{field content},\text{gauge},\text{basis},\text{Fock set},
 \text{regulators},\text{canonical terms},\text{induced basis},\\
 \text{counterterms},\text{renormalization conditions},\text{version}
 \end{gathered}
 \right).
 \label{eq:hamiltonian-identity}
\end{equation}
Two matrices with the same dimension but different \(\mathfrak h_r\) are
different theories and may not be compared as though only numerical resolution
had changed.
\end{definition}

\subsection{Free invariant mass}

For a sector \(\nu\) with partons of momentum fractions \(x_i>0\), intrinsic
transverse momenta \(\bm\kappa_{iT}\), and masses \(m_i^{(r)}\),
\begin{equation}
 \cM_{0,\nu}^{2}
 =\sum_{i\in\nu}
 \frac{\bm\kappa_{iT}^{2}+m_i^{2}}{x_i},
 \qquad
 \sum_i x_i=1,
 \qquad
 \sum_i\bm\kappa_{iT}=0.
 \label{eq:free-invariant-mass}
\end{equation}
The masses in \cref{eq:free-invariant-mass} are scheme- and sector-indexed bare
or effective parameters.  They are not identified with \(\overline{\rm MS}\)
current masses without a matching calculation.

\subsection{Canonical interaction terms}

In light-front gauge \(A^+=0\), the retained canonical operator set includes
quark--gluon emission and absorption, following the standard canonical
light-front QCD reduction \cite{BrodskyPauliPinsky1998},
\begin{equation}
 V_{qg}^{\can}
 =g_s\int [\dd x]\,
 \bar\psi\gamma^\mu t^a\psi A_\mu^a,
 \label{eq:qg-vertex}
\end{equation}
three- and four-gluon vertices where the sector content permits them, and the
instantaneous interactions generated by eliminating constrained field
components.  Schematically,
\begin{align}
 V_{qq}^{\inst}
 &\sim g_s^2
 \int [\dd x]\,
 \bigl(\bar\psi\gamma^+t^a\psi\bigr)
 \frac{1}{(i\partial^+)^2}
 \bigl(\bar\psi\gamma^+t^a\psi\bigr),
 \label{eq:instantaneous-quark}\\
 V_{qgg}^{\inst}
 &\sim g_s^2
 \int [\dd x]\,
 \bar\psi\gamma^\mu t^a A^a_\mu
 \frac{\gamma^+}{i\partial^+}
 \gamma^\nu t^b A^b_\nu\psi .
 \label{eq:instantaneous-fermion}
\end{align}
Their kernels, endpoint prescriptions, color factors, and Hermitian-conjugate
blocks are part of the Hamiltonian identity.

\begin{requirement}[Canonical-term closure]
Every canonical vertex included in one direction between retained sectors must
have its Hermitian-conjugate block, all instantaneous partners required by the
selected gauge reduction, and the counterterm structures required by the
renormalization conditions.  Omitting a partner is a declared approximation
with a Ward-identity defect; it is not a silent simplification.
\label{req:canonical-closure}
\end{requirement}

\subsection{Block structure}

With sector projectors \(P_3,P_{4g},P_{5q}\) for
\(qqq\), \(qqqg\), and \(qqqq\bar q\), respectively,
\begin{equation}
 \cM_r^2=
 \begin{pmatrix}
  H_{33} & V_{3,4g} & V_{3,5q}\\
  V_{4g,3} & H_{4g,4g} & V_{4g,5q}\\
  V_{5q,3} & V_{5q,4g} & H_{5q,5q}
 \end{pmatrix}_r .
 \label{eq:block-hamiltonian}
\end{equation}
At the first implementation level, canonical one-gluon emission produces
\(V_{3,4g}\), while pair creation and annihilation may connect the retained
sectors directly or through induced operators depending on the chosen
Hamiltonian reduction.  A zero block must mean either a forbidden quantum
number or an explicitly declared truncation.

\begin{proposition}[Hermitian block criterion]
If each diagonal block is self-adjoint on the finite regulated domain and
\(V_{\nu\mu}=V_{\mu\nu}^{\dagger}\), then
\(\cM_r^2=(\cM_r^2)^\dagger\).  This algebraic condition is necessary but not
sufficient for physical renormalization, gauge consistency, or continuum
convergence.
\label{prop:block-hermiticity}
\end{proposition}

\section{Induced confinement, chiral dynamics, and omitted sectors}
\label{sec:induced}

\subsection{Confinement as a resolution-dependent induced operator}

A finite low-resolution partonic Hamiltonian does not contain the full soft and
higher-sector dynamics that produces confinement in continuum QCD.  The use of
effective interactions and Hamiltonian similarity flows in truncated spaces is
well established, but their coefficients and operator content must remain
resolution dependent \cite{Wilson1994,GlazekWilson1993,Wegner1994}.  The first
state calculation therefore admits an induced operator
\begin{equation}
 V_{\conf,r}^{\ind}
 =\kappa_{T,r}^{4}\,\cO_{T}^{\conf}
 +\kappa_{L,r}^{4}\,\cO_{L}^{\conf}
 +\sum_{n>1}d_{n,r}^{\conf}\cO_{n}^{\conf},
 \label{eq:induced-confinement}
\end{equation}
where the operator basis respects color, flavor, \(\Jz\), parity, permutation
symmetry, and cluster limits.  The displayed transverse and longitudinal
operators may be represented in Jacobi variables by a harmonic or another
smooth infrared form at the coarsest truncations, but their coefficients are
resolution dependent.

\begin{decision}[Confinement ablation]
Every accepted truncation tower contains both:
\begin{enumerate}
 \item an induced-confinement trajectory in which
 \(V_{\conf,r}^{\ind}\) is refitted at each resolution; and
 \item a canonical-QCD trajectory in which its coefficient is driven toward
 zero as additional basis and Fock content is included.
\end{enumerate}
A confining term is accepted as a useful low-resolution operator only if
matched observables stabilize and its coefficient exhibits an interpretable
flow.  It may not be treated as a universal immutable potential.
\label{dec:confinement-ablation}
\end{decision}

The induced operator is not allowed to duplicate one-gluon exchange or an
explicit pionlike five-parton channel.  Its provenance record therefore lists
which infrared regions and excluded sectors it represents.

\subsection{Chiral sea operator}

An explicit \(qqqq\bar q\) sector is necessary but not sufficient to generate
the observed flavor structure of the sea.  Light-front Hamiltonian EFT studies
with explicit multi-pion sectors illustrate how nonperturbative chiral Fock
components can organize this physics \cite{Cao2026}.  At low resolution the Hamiltonian
must support strong pseudoscalar-isovector correlations.  We write the leading
induced chiral block schematically as
\begin{equation}
 V_{\chi,r}^{\ind}
 =G_{\chi,r}
 \sum_{A=1}^{3}
 \int [\dd\Gamma]
 \left(\bar\psi i\gamma_5\tau^A\psi\right)
 \cF_{\chi,r}
 \left(\bar\psi i\gamma_5\tau^A\psi\right)
 +\cdots,
 \label{eq:chiral-induced}
\end{equation}
where \(\cF_{\chi,r}\) is a regulated kernel and the ellipsis denotes the
scalar or derivative partners required by the chosen chiral realization.
Equation~\eqref{eq:chiral-induced} is an operator template, not a commitment to
one local four-fermion approximation.  Its admissible implementation must be
calibrated in the meson and nucleon sectors and must reproduce the appropriate
chiral limits to the declared order.

An alternative but matched representation introduces an explicit effective
pion channel at the hadronic resolution.  The partonic five-body and hadronic
pion descriptions are then related by a matching map and overlap subtraction,
\begin{equation}
 \cO_{5q}^{\rm matched}
 =\cO_{5q}^{\rm partonic}
 +\cO_{N\pi}^{\rm hadronic}
 -\cO_{\rm overlap}^{\chi}.
 \label{eq:chiral-overlap-subtraction}
\end{equation}
They are not independent additive mechanisms.

\begin{requirement}[Sea flavor is dynamical]
The proton's \(\bar d-\bar u\) structure must arise from the shared Hamiltonian,
its chiral operators, and its five-parton amplitudes.  It may not be produced by
a large ad hoc \(u-d\) constituent-mass split or by assigning independent
antiquark distributions after diagonalization.
\label{req:sea-dynamical}
\end{requirement}

\subsection{Spin-restoration operators}

Truncating the basis and Fock space breaks dynamical transverse rotations.  The
induced basis may therefore contain symmetry-restoration operators
\begin{equation}
 V_{\rm spin,r}^{\ind}
 =\sum_a c_{a,r}^{\rm rot}\cO_{a}^{\rm rot}
 +\sum_b c_{b,r}^{\rm hf}\cO_{b}^{\rm helicity\ flip},
 \label{eq:spin-restoration}
\end{equation}
provided each term has a declared origin, power counting, and vanishing or
stable continuum behavior.  Such operators are constrained by masses,
current angular conditions, Pauli--Lubanski diagnostics, and multiplet
relations.  They cannot be tuned independently for every helicity amplitude.

\subsection{Omitted-sector remainder}

Let \(P_r\) project onto the retained space and \(Q_r=1-P_r\).  Formally, the
exact energy-dependent effective Hamiltonian is
\begin{equation}
 H_{\eff,r}(E)
 =P_rHP_r
 +P_rHQ_r\frac{1}{E-Q_rHQ_r}Q_rHP_r.
 \label{eq:feshbach-hamiltonian}
\end{equation}
The induced terms in \cref{eq:master-hamiltonian} approximate the second term.
Their error is recorded as an operator-valued remainder rather than one scalar
band,
\begin{equation}
 \Delta\cM_{r}^{2,\trunc}
 =\sum_{A\notin\cB_r}
 \delta c_{A,r}\cO_A,
 \qquad
 \delta c_{A,r}\sim Q_r^{\nu_A},
 \label{eq:operator-remainder}
\end{equation}
where \(Q_r\) denotes one or more expansion parameters associated with the
basis, Fock, or resolution truncation.

\section{Fock hierarchy and sector content}
\label{sec:fock}

\subsection{Directed sector tower}

The microscopic program uses a nested physical hierarchy even when the
underlying basis spaces are related by comparison rather than literal inclusion:
\begin{align}
 \cF^{(0)}&=\{qqq\},
 \label{eq:fock0}\\
 \cF^{(1)}&=\{qqq,qqqg\},
 \label{eq:fock1}\\
 \cF^{(2)}&=\{qqq,qqqg,qqqu\bar u,qqqd\bar d\},
 \label{eq:fock2}\\
 \cF^{(3)}&=\cF^{(2)}\cup\{qqqgg,qqqq\bar q g,s\bar s,\ldots\}.
 \label{eq:fock3}
\end{align}
Each level is solved and renormalized independently.  A sector probability may
change substantially between levels without implying a physical contradiction;
matched masses, currents, and distributions must stabilize.

\subsection{Tier 0: valence benchmark}

The \(qqq\) calculation establishes:
\begin{itemize}
 \item exact color-singlet and permutation bases;
 \item center-of-mass factorization;
 \item mass and current normalization;
 \item \(\Jz\) and \(\Lz\) bookkeeping;
 \item the baseline for induced-confinement and rotational-restoration flows.
\end{itemize}
It cannot generate antiquarks, dynamical gluons, genuine quark--gluon
correlations, or physical naive-\(T\)-odd phases.  It is therefore a benchmark,
not the first admissible common state bundle.

\subsection{Tier 1: one dynamical gluon}

The \(qqqg\) sector introduces:
\begin{itemize}
 \item gluon PDFs, helicity, OAM, and field-strength operator overlaps;
 \item quark self-energy and vertex dressing;
 \item the intermediate state required by the first Wilson-line cut pilot;
 \item color-octet three-quark subspaces coupled to an adjoint gluon;
 \item genuine quark--gluon and twist-three operator support.
\end{itemize}
The color basis contains every independent singlet multiplicity.  A singlet
\(qqq\) tensor multiplied by a free gluon is forbidden.

\subsection{Tier 2: explicit light sea}

The five-parton sector contains all independent color-singlet multiplicities,
full antisymmetry among identical quarks, and explicit active quark and
antiquark slots.  The state is written
\begin{equation}
 |N,\Lambda\rangle_{5q}
 =\sum_{q=u,d}
 \sum_{\alpha_{5q}}
 \int [\dd\Gamma_5]\,
 \psi^{\Lambda}_{5q,\alpha_{5q}}
 |qqqq\bar q;\alpha_{5q}\rangle.
 \label{eq:five-quark-state}
\end{equation}
The basis must be capable of representing both meson--baryon-like clusters and
non-cluster color recouplings.  Cluster language is an analysis basis, not a
probability decomposition unless the channels are orthogonal.

\subsection{Tier 3 and closure sectors}

The sectors \(qqqgg\) and \(qqqq\bar qg\) are expected to become necessary for
higher Wilson orders, gluon self-interactions, complete Ward closure, and
renormalization stability.  Strange pairs are added when the target
observables, matching scale, or threshold treatment requires them.

\begin{requirement}[Sector-promotion rule]
A missing sector is promoted from remainder to explicit state space when at
least one of the following occurs:
\begin{enumerate}
 \item a Ward or symmetry defect does not decrease under the current tower;
 \item an induced coefficient becomes unnatural or strongly regulator
 dependent;
 \item a withheld observable is sensitive to a quantum number unavailable in
 the retained sectors;
 \item the estimated omitted-sector remainder exceeds the accepted precision;
 \item an operator insertion has no closed intermediate-state support.
\end{enumerate}
\label{req:sector-promotion}
\end{requirement}

\section{Symmetry-adapted BLFQ basis}
\label{sec:basis}

\subsection{Single-particle labels}

A single-particle basis state is labeled by
\begin{equation}
 \alpha_i=
 \left(
 a_i,k_i,n_i,m_i,\lambda_i,f_i,c_i
 \right),
 \label{eq:single-particle-label}
\end{equation}
where \(a_i\) is the parton species, \(k_i>0\) is the longitudinal mode,
\((n_i,m_i)\) are transverse two-dimensional harmonic-oscillator labels,
\(\lambda_i\) is light-front helicity, \(f_i\) is flavor, and \(c_i\) is a
color basis label.  This is the standard BLFQ organization adapted here to a
fully symmetry-projected multi-sector state \cite{Vary2010,Zhao2014,Mondal2025}.

The total longitudinal resolution is
\begin{equation}
 \sum_i k_i=K,
 \qquad
 x_i=\frac{k_i}{K},
 \label{eq:longitudinal-K}
\end{equation}
and the transverse truncation is
\begin{equation}
 \sum_i\left(2n_i+|m_i|+1\right)\leq N_{\max}.
 \label{eq:nmax}
\end{equation}
The oscillator scale \(b\) is a basis parameter, not a physical transverse
width.

\subsection{Longitudinal boundary conditions}

The initial discretization uses antiperiodic longitudinal modes for quarks and
antiquarks and periodic nonzero modes for gluons.  The physical output records
this choice.  The corresponding zero-mode policy is not simply ``omit the
zero mode'': it is
\begin{equation}
 \mathfrak z_r=
 \left(
 \begin{gathered}
 \text{excluded modes},\text{constraint equations},
 \text{induced zero-mode operators},\\
 \text{endpoint regulator},\text{closure tests}
 \end{gathered}
 \right).
 \label{eq:zero-mode-policy}
\end{equation}
A later zero-mode-resolved tower may replace this policy.

\subsection{Ultraviolet and infrared interpretation}

For the oscillator basis,
\begin{equation}
 \Lambda_{\IR}^{T}\sim\frac{b}{\sqrt{N_{\max}}},
 \qquad
 \Lambda_{\UV}^{T}\sim b\sqrt{N_{\max}},
 \label{eq:ho-cutoffs}
\end{equation}
up to sector- and observable-dependent constants.  The longitudinal
resolution and endpoint cutoff are governed by \(K\) and \(\epsilon_x\).
Changing \(b\) at fixed \(N_{\max}\) moves both cutoffs; it is therefore a
regulator/basis variation, not a profile uncertainty.

\subsection{Intrinsic and center-of-mass factorization}

The many-body basis is transformed to intrinsic Jacobi coordinates or equipped
with a Lawson term,
\begin{equation}
 \cM_{r,\beta}^{2}
 =\cM_r^2
 +\beta_{\COM}
 \left(
 H_{\COM}-E_{\COM,0}
 \right),
 \label{eq:lawson-term}
\end{equation}
so spurious center-of-mass excitations are separated from the intrinsic
spectrum.  Physical eigenvalues and matrix elements must be stable over an
accepted range of \(\beta_{\COM}\).

\begin{proposition}[Center-of-mass gate]
A state is admissible for partonic matrix elements only if its center-of-mass
contamination, measured through the Lawson excitation and direct factorization
residual, is below the declared tolerance.  Orthogonality of basis vectors
alone does not establish intrinsic factorization.
\label{prop:com-gate}
\end{proposition}

\subsection{Total angular momentum projection}

The kinematical generator satisfies
\begin{equation}
 \Jz=\sum_i\left(m_i+\lambda_i\right),
 \label{eq:jz-basis}
\end{equation}
with the appropriate helicity convention for antiquarks and gluons.  The
Hamiltonian is assembled separately in each \(\Jz\), baryon-number, charge,
flavor, parity-class, and color-singlet block.  Transverse rotations remain
dynamical and are monitored through the diagnostics of
\cref{sec:poincare}.

\section{Color, statistics, and physical-state projection}
\label{sec:color-statistics}

\subsection{Exact color singlets}

For every sector, the color basis is the full invariant subspace
\begin{equation}
 \Hil_{\nu,\mathrm{color}=1}
 =\operatorname{Inv}_{SU(3)}
 \left[
 \bigotimes_{i\in\nu}R_{c_i}
 \right].
 \label{eq:color-invariant-space}
\end{equation}
Independent singlet multiplicities are retained as basis labels.  Each color
tensor \(C_\alpha\) satisfies
\begin{equation}
 \left(\sum_i T_i^A\right)C_\alpha=0
 \qquad \forall A=1,\ldots,8.
 \label{eq:color-generator-test}
\end{equation}
For antiquarks \(T_{\bar3}^A=-(t^A)^{\mathsf T}\), and for gluons
\((T_8^A)_{bc}=-if^{Abc}\).

\subsection{Full identical-fermion antisymmetry}

The many-quark basis must satisfy
\begin{equation}
 \cP_{ij}|\Psi\rangle=-|\Psi\rangle
 \label{eq:fermion-antisymmetry}
\end{equation}
for exchange of any two identical quarks, including exchanges across apparent
cluster partitions.  It is insufficient to choose an ordered flavor--spin
configuration and one global singlet tensor without projecting the complete
state.

A practical construction uses irreducible representations of the relevant
permutation group, with color, flavor, spin, orbital, and longitudinal labels
coupled to the totally antisymmetric representation.

\begin{requirement}[Antisymmetry before diagonalization]
Full identical-fermion antisymmetry is imposed on the basis before Hamiltonian
assembly.  Post hoc symmetrization of selected eigenvectors is not accepted,
because it can alter matrix dimensions, vertex multiplicities, and current
normalizations.
\label{req:antisymmetry}
\end{requirement}

\subsection{Physical projector}

The regulated physical sector is
\begin{equation}
 \Hil_{\nu,\phys}^{(r)}
 =P_{\rm color=1}
 P_{\rm stat}
 P_{B,Q,I_3,\Jz,\Pi}
 P_{\rm gauge}^{(r)}
 \cK_{\nu}^{(r)}.
 \label{eq:physical-projector}
\end{equation}
The residual gauge projector encodes the constraints and boundary data left
after gauge fixing.  Color-singlet projection is necessary but does not by
itself prove gauge consistency.

\section{Regulator architecture}
\label{sec:regulators}

\subsection{Vertex invariant-mass regulator}

In addition to the basis cutoffs, sector-changing vertices use a common smooth
invariant-mass-change regulator,
\begin{equation}
 f_{\lambda}
 \left(\Delta\cM^2\right)
 =\exp\left[-\frac{(\Delta\cM^2)^2}{\lambda^4}\right],
 \label{eq:vertex-regulator}
\end{equation}
or a declared alternative with the same symmetry and limiting properties.
Here
\(\Delta\cM^2=\cM_{\mathrm{out},0}^2-\cM_{\mathrm{in},0}^2\).  The same regulator
identity is used in Hermitian-conjugate blocks and associated currents.

The parameter \(\lambda\) is a Hamiltonian resolution, not the TMD rapidity or
Collins--Soper scale.  Its notation and type remain distinct from \(\mu\) and
\(\zeta\).

\subsection{\texorpdfstring{Small-\(x\)}{Small-x} and endpoint treatment}

The basis itself enforces \(x_i\geq 1/(2K)\) or \(1/K\), depending on boundary
conditions.  If additional endpoint damping is required, it is represented by
a typed regulator
\begin{equation}
 R_{\epsilon_x}(\{x_i\})
 =\prod_i r_{a_i}(x_i;\epsilon_x)
 \label{eq:endpoint-regulator}
\end{equation}
and accompanied by an endpoint counterterm or remainder analysis.  The
regulator cannot be mistaken for a predicted small-\(x\) power.

\subsection{Instantaneous-kernel prescription}

The inverse derivatives in instantaneous interactions require a declared
principal-value, zero-mode-subtracted, or finite-volume prescription.  The
prescription is included in \(\mathfrak h_r\) and shared by the Hamiltonian and
current operators.

\subsection{Regulator covariance}

The regulator set must preserve the exact kinematical symmetries of the finite
basis and may violate dynamical symmetries only in a way monitored by the
restoration program.  A regulator that depends on one TMD projection or one
active flavor is prohibited unless it represents a physically distinct
operator threshold.

\begin{requirement}[No regulator migration]
Basis parameters \((K,N_{\max},b)\), Hamiltonian resolution \(\lambda\),
endpoint regulator \(\epsilon_x\), ultraviolet scheme, TMD rapidity regulator,
and numerical quadrature cutoffs are different typed quantities.  No API may
supply one in place of another by a scalar alias.
\label{req:no-regulator-migration}
\end{requirement}

\section{Renormalization datum and parameter ownership}
\label{sec:renormalization}

\subsection{Complete renormalization datum}

At every truncation level, the model carries
\begin{equation}
 \mathfrak R_r=
 \left(
 \cR_r,
 \theta_{\bare}^{(r)},
 \delta\theta_{\ct}^{(r)},
 \theta_{\ind}^{(r)},
 \{\cC_i\}_{i\in\calib},
 \cS_r,
 \cZ_r,
 \Delta_r
 \right),
 \label{eq:renormalization-datum}
\end{equation}
where \(\cR_r\) is the regulator set, \(\theta_{\bare}^{(r)}\) are bare
couplings and masses, \(\delta\theta_{\ct}^{(r)}\) are counterterms,
\(\theta_{\ind}^{(r)}\) are induced-operator coefficients,
\(\{\cC_i\}\) are renormalization conditions, \(\cS_r\) is the scheme,
\(\cZ_r\) is the operator-renormalization map, and \(\Delta_r\) is the
truncation discrepancy model.

\begin{principle}[Renormalization is a trajectory]
A parameter set at one \((K,N_{\max},b,\lambda,\cF)\) is not the model.  The
model is the family of refitted parameters and matched observables along a
directed truncation trajectory.  Convergence is assessed after imposing the
same physical renormalization conditions at every point.
\label{pr:renorm-trajectory}
\end{principle}

\subsection{Parameter ownership}

The initial shared parameter blocks are:
\begin{longtable}{L{0.22\textwidth}L{0.28\textwidth}L{0.42\textwidth}}
\toprule
Block & Representative parameters & Ownership and restrictions\\
\midrule
\endhead
Kinetic/QCD & \(m_{u}^{\bare},m_d^{\bare},m_g^{\bare},g_s^{\bare}\) &
Canonical kinetic and gauge vertices.  Strong isospin symmetry uses one light
mass at the reference point; \(m_d-m_u\) and QED are controlled breaking
operators.\\
Mass counterterms & \(\delta m_{q,\nu}^{2},\delta m_{g,\nu}^{2}\) &
Fock-sector- and truncation-indexed self-energy renormalization.  They are
determined by shared pole or invariant-mass conditions, not fitted to TMD
shapes.\\
Vertex counterterms & \(\delta g_{\nu\mu}^{(r)}\), helicity-flip terms &
Constrained by Ward identities, vertex matrix elements, and symmetry-restoration
conditions.\\
Induced confinement & \(\kappa_{T,r},\kappa_{L,r},d_{n,r}^{\conf}\) &
Low-resolution omitted-sector operators; refitted along the tower and tested
against a zero-confinement trajectory.\\
Chiral sector & \(G_{\chi,r}\), range and derivative coefficients &
Shared meson/nucleon operator basis.  Calibrated to chiral and sea-sensitive
observables, with explicit overlap subtraction if a hadronic pion channel is
used.\\
Current/operator & \(Z_J^{(r)},c_{J,a}^{(r)},Z_{\GTMD}^{(r)}\) &
Consistent one-, two-, instantaneous-, and induced-current terms.  Cannot be
fitted independently for each form factor or named GTMD.\\
Discrepancy & \(\delta c_{A,r}\) with power-counting priors &
Explicit omitted-operator basis.  May not duplicate an included sector or
absorb arbitrary residual functions.\\
\bottomrule
\end{longtable}

\subsection{Strong isospin and controlled breaking}

The reference strong Hamiltonian satisfies
\begin{equation}
 m_u^{\bare}=m_d^{\bare}=m_\ell^{\bare},
 \qquad
 [\cM_{\mathrm{strong}}^{2},\bm I]=0.
 \label{eq:isospin-reference}
\end{equation}
Charge-symmetry breaking is added through
\begin{equation}
 \delta\cM_{\rm CSB}^{2}
 =\delta m_{d-u}\,\cO_{d-u}
 +\alpha_{\rm em}\cO_{\rm QED}
 +\cdots.
 \label{eq:csb-operators}
\end{equation}
The proton and neutron are correlated members of the same isospin multiplet.
Their flavor differences emerge from the state and controlled breaking, not
from independent wave-function fits.

\begin{requirement}[Natural parameter count]
The number of freely calibrated microscopic parameters must remain
substantially smaller than the number of observable families predicted.
Counterterms required to impose a renormalization condition are recorded
separately from physical low-energy constants and do not count as independent
named-observable normalizations.
\label{req:natural-parameter-count}
\end{requirement}

\section{Fock-sector-dependent renormalization}
\label{sec:fsdr}

\subsection{Why sector dependence appears}

At finite Fock truncation, a constituent in the highest retained sector cannot
fluctuate into the same omitted states as a constituent in a lower sector.
Fock-sector-dependent renormalization provides a systematic way to encode this
asymmetry and preserve Ward constraints in controlled truncations
\cite{Karmanov2005,Karmanov2008,Karmanov2012}.
Therefore its self-energy and vertex dressing differ.  A common bare mass or
coupling inserted blindly in every block does not generally satisfy the same
renormalization conditions.

We write sector-indexed parameters as
\begin{equation}
 m_{a,\nu}^{2}
 =m_{a,\phys}^{2}+\delta m_{a,\nu}^{2},
 \qquad
 g_{\nu\mu}
 =g_{\phys}+\delta g_{\nu\mu},
 \label{eq:fsdr-parameters}
\end{equation}
with consistency relations that reduce the number of independent quantities.

\subsection{Mass renormalization conditions}

The nucleon pole condition is
\begin{equation}
 \det\left[
 \cM_r^{2}(\theta_r)-M_N^2\bbone
 \right]=0.
 \label{eq:nucleon-pole}
\end{equation}
This condition alone is insufficient to determine all sector counterterms.
Additional conditions may include:
\begin{enumerate}
 \item constituent two-point conditions in controlled one-body or
 quark--gluon subproblems;
 \item equality of selected self-energy insertions under spectator changes;
 \item matching to meson-sector masses for shared quark/chiral parameters;
 \item stability of the nucleon pole under sector promotion.
\end{enumerate}

The highest-sector mass treatment must not switch silently from effective to
current masses.  Any sector-dependent mass is part of \(\mathfrak R_r\) and
must be justified by an explicit renormalization condition or induced operator.

\subsection{Vertex and Ward renormalization}

For a conserved vector current, the charge condition is
\begin{equation}
 F_1^p(0)=1,
 \qquad
 F_1^n(0)=0,
 \label{eq:charge-renorm}
\end{equation}
and the underlying vertex and wave-function factors satisfy a regulated
Ward--Takahashi relation.  In schematic notation,
\begin{equation}
 q_\mu\Gamma_r^\mu(p+q,p)
 =S_r^{-1}(p+q)-S_r^{-1}(p)
 +\Delta_{\WTI,r},
 \label{eq:wti}
\end{equation}
where \(\Delta_{\WTI,r}\) is a measured truncation and regulator defect.
The renormalization program must reduce this defect along the tower.

\subsection{Sector-independent observables}

The purpose of \FSDR is not to make sector probabilities equal.  It is to make
matched observables insensitive to the arbitrary location of the Fock cutoff:
\begin{equation}
 \cO_i^{(r')}
 -\cO_i^{(r)}
 =\cO\!\left(Q_{r'r}^{\nu_i}\right)
 \label{eq:observable-convergence}
\end{equation}
for renormalization and withheld observables after the required comparison
maps have been applied.

\begin{requirement}[FSDR closure]
The first microscopic state bundle is not accepted unless the mass, charge,
selected current matrix elements, and common-parent reduction moments remain
stable under at least one basis enlargement and one Fock-sector enlargement
with counterterms refitted.  A single successful diagonalization does not
satisfy this requirement.
\label{req:fsdr-closure}
\end{requirement}

\section{Consistent currents and partonic operators}
\label{sec:operators}

\subsection{Hamiltonian and operator must transform together}

If a unitary, similarity, or Feshbach transformation \(U_r\) changes the
Hamiltonian representation \cite{Feshbach1958,Okubo1954,LeeSuzuki1980,GlazekWilson1993},
\begin{equation}
 H_r\mapsto H_r'=U_rH_rU_r^\dagger,
 \label{eq:ham-transform}
\end{equation}
then every observable operator changes as
\begin{equation}
 \cO_r\mapsto\cO_r'=U_r\cO_rU_r^\dagger.
 \label{eq:operator-transform}
\end{equation}
Eliminating a sector from \(H\) while replacing the current by \(P\cO P\)
misses induced contributions.  For a Feshbach wave operator
\begin{equation}
 \omega_r(E)=\frac{1}{E-Q_rH_rQ_r}Q_rH_rP_r,
 \label{eq:wave-operator}
\end{equation}
the corresponding effective operator is
\begin{equation}
 \cO_{\eff,r}(E',E)
 =P_r\left[1+\omega_r^\dagger(E')\right]
 \cO
 \left[1+\omega_r(E)\right]P_r,
 \label{eq:feshbach-operator}
\end{equation}
with an associated norm kernel.

\begin{theorem}[Representation invariance of complete matrix elements]
If the state, Hamiltonian, norm kernel, and operator are transformed
consistently, the full-space and effective-space matrix elements agree within
the declared truncation accuracy.  A decomposition into ``one-body off-shell''
and ``two-body induced'' contributions is representation dependent, while the
sum is physical.
\label{thm:representation-invariance}
\end{theorem}

\subsection{Electromagnetic current}

The regulated current is
\begin{equation}
 J_r^\mu
 =J_{(1),r}^\mu
 +J_{\inst,r}^\mu
 +J_{(2),r}^\mu
 +J_{\ind,r}^\mu
 +\delta J_{\ct,r}^\mu.
 \label{eq:em-current}
\end{equation}
It is constructed by gauging the Hamiltonian and regulator or by an equivalent
operator-matching procedure.  Form factors evaluated with different current
components must agree up to measured rotational and truncation defects.

\subsection{Axial and pseudoscalar operators}

The axial current and pseudoscalar density obey a regulated PCAC relation,
\begin{equation}
 \partial_\mu A_r^{\mu,A}
 =2m_{\ell,r}P_r^A
 +f_{\pi,r}m_{\pi,r}^2\phi_r^A
 +\Delta_{\PCAC,r}^A.
 \label{eq:pcac}
\end{equation}
The chiral induced interaction and axial current must share parameters and
counterterms.  A chiral sea operator that improves \(\bar d-\bar u\) but
produces an uncontrolled \(\Delta_{\PCAC}\) is not acceptable.

\subsection{Energy--momentum tensor}

The quark and gluon EMT operators are matched consistently to the Hamiltonian,
so that
\begin{equation}
 \sum_{a=q,\bar q,g}
 \avg{x}_a^{(r)}
 =1+\Delta_{P^+,r},
 \label{eq:momentum-sum}
\end{equation}
and the angular-momentum decomposition is evaluated in a declared gauge and
scheme,
\begin{equation}
 \frac12
 =\frac12\Delta\Sigma_r
 +L_{q,r}
 +\Delta G_r
 +L_{g,r}
 +\Delta_{J,r}.
 \label{eq:spin-sum}
\end{equation}
The residuals are separately reported; they are not absorbed into one
normalization factor.

\subsection{GTMD operator bundle}

The common partonic operator set includes the complete decorated identities
needed by Volume~II.  At the state level, the first required exports are:
\begin{itemize}
 \item all nucleon helicity matrix elements for quark, antiquark, and gluon
 leading-twist projectors at zero Wilson order;
 \item nonzero-\(\DT\) momentum-fiber support and exact recoil metadata;
 \item straight-link and staple-ready operator identities;
 \item local-current and EMT reductions from the same operator basis;
 \item derivative and phase conventions required by OAM projections.
\end{itemize}

\begin{requirement}[Closed operator basis]
Before an LF-to-QCD matching calculation begins, the regulated Hamiltonian
operator basis must close under the retained symmetries, sector elimination,
and renormalization to the declared order.  ``One current per observable'' is
not a closed basis.
\label{req:closed-operator-basis}
\end{requirement}

\section{Poincar\'e and rotational restoration}
\label{sec:poincare}

\subsection{Exact and dynamical generators}

Light-front kinematics preserves seven kinematical generators, including
\(P^+\), \(\bm P_T\), and \(J^z\), while transverse rotations are dynamical.
The finite basis therefore enforces exact \(\Jz\) but measures the remaining
algebraic defects,
\begin{equation}
 \Delta_{ab}^{(r)}
 =\left\|
 [G_a^{(r)},G_b^{(r)}]
 -if_{ab}^{\ \ c}G_c^{(r)}
 \right\|_{\cS_{\rm test}}.
 \label{eq:poincare-defect}
\end{equation}
The norm is evaluated on a declared low-energy test subspace, not over the
entire cutoff matrix where ultraviolet artifacts may dominate.

\subsection{Pauli--Lubanski diagnostics}

The Pauli--Lubanski vector
\begin{equation}
 W^\mu=-\frac12\epsilon^{\mu\nu\rho\sigma}
 J_{\nu\rho}P_\sigma
 \label{eq:pauli-lubanski}
\end{equation}
obeys
\begin{equation}
 W^2=-M^2s(s+1)
 \label{eq:w2}
\end{equation}
for an exact massive spin-\(s\) state.  The finite-basis residual and the
consistency of different current helicity extractions are part of the
rotational-restoration report.

\subsection{Angular conditions}

For spin-one-half nucleons, rotational covariance is tested through equality of
form factors extracted from independent current components and frames.  For
the later deuteron state, the spin-1 angular condition becomes an additional
gate.  Volume~VII requires the nucleon current representation to carry the
metadata and covariance needed for that later test.

\subsection{Restoration strategy}

Restoration proceeds in this order:
\begin{enumerate}
 \item enlarge \(K\), \(N_{\max}\), and the Fock space;
 \item refit required counterterms under the same conditions;
 \item add only symmetry-allowed induced operators supported by the defect
 pattern;
 \item verify that their coefficients stabilize or decrease;
 \item withhold independent current and GTMD relations as tests.
\end{enumerate}
A direct fit of every current component to force equality is prohibited.

\section{Sparse eigensolver and state tracking}
\label{sec:solver}

\subsection{Matrix representation}

The Hamiltonian is block sparse in exact quantum numbers and color/permutation
multiplicity.  The production representation should be matrix free whenever
possible:
\begin{equation}
 \bm y=\cM_r^2\bm x
 \label{eq:matrix-free}
\end{equation}
without materializing all zero blocks.  Vertex kernels, recoupling tensors, and
regulators are cached as typed objects.

\subsection{Eigensolver hierarchy}

The solver stack is:
\begin{enumerate}
 \item exact diagonalization for small analytic and low-dimensional blocks;
 \item Lanczos or implicitly restarted Arnoldi for isolated low states;
 \item block Davidson or Jacobi--Davidson for near-degenerate proton/neutron
 and parity/helicity multiplets;
 \item shift-invert or contour methods only with a documented preconditioner
 and spectral target.
\end{enumerate}
Every eigenpair carries the residual
\begin{equation}
 r_n=\frac{\norm{\cM_r^2\psi_n-M_{n,r}^2\psi_n}}
 {\norm{\cM_r^2}\norm{\psi_n}+|M_{n,r}^2|\norm{\psi_n}}.
 \label{eq:eigen-residual}
\end{equation}

\subsection{State identity across truncations}

The nucleon state at \(r\) is matched to states at \(r'\) using quantum
numbers, comparison-map overlap, current fingerprints, and sector/OAM content.
For a nondegenerate state,
\begin{equation}
 \cO_{rr'}
 =\abs{\inner{\Psi_N^{(r')}}{I_{r\to r'}\Psi_N^{(r)}}}^2.
 \label{eq:state-overlap}
\end{equation}
Within a nearly degenerate subspace, principal angles and subspace projectors
replace individual-vector overlaps.  Phase conventions are fixed by a declared
reference component and transported continuously.

\begin{requirement}[No eigenvalue-only tracking]
A state may not be identified across truncations only because it is the lowest
eigenvalue.  Avoided crossings and basis artifacts require overlap, current,
and quantum-number diagnostics.
\label{req:no-eigenvalue-tracking}
\end{requirement}

\subsection{Differentiation}

For a normalized isolated eigenstate,
\begin{equation}
 \frac{\partial M_N^2}{\partial\theta_a}
 =\left\langle\Psi_N\left|
 \frac{\partial\cM^2}{\partial\theta_a}
 \right|\Psi_N\right\rangle
 \label{eq:hellmann-feynman}
\end{equation}
provides the Hellmann--Feynman derivative.  State derivatives and observable
Jacobians use adjoint or implicit differentiation with degeneracy safeguards.
Finite differences remain an independent validation route.

\section{Calibration and withheld validation}
\label{sec:calibration}

\subsection{Calibration philosophy}

The first Hamiltonian calibration must be small enough that the state predicts
more than it absorbs.  The initial recommended calibration set is:
\begin{enumerate}
 \item proton and neutron masses in the isospin-symmetric/controlled-CSB
 organization;
 \item proton charge normalization and one selected magnetic or radius
 combination;
 \item pion mass and decay/axial condition for the shared chiral operator, when
 that operator is active;
 \item one selected axial or tensor matrix element;
 \item renormalization-subproblem conditions needed to fix sector counterterms.
\end{enumerate}

The following families should remain withheld at the first pass:
\begin{itemize}
 \item additional electromagnetic form factors and radii;
 \item neutron magnetic and charge structure not used in calibration;
 \item the other of axial or tensor charge;
 \item flavor-resolved PDF moments and shapes;
 \item gluon momentum and helicity fractions;
 \item sea asymmetry beyond any one calibration moment;
 \item OAM-sensitive GTMD relations;
 \item all TMD-specific shapes and naive-\(T\)-odd observables.
\end{itemize}

\subsection{No model-scale fitting by evolution alone}

A low model scale cannot be chosen solely so that DGLAP evolution reproduces a
selected PDF moment.  The Hamiltonian resolution, LF-to-QCD matching, and
collinear evolution are distinct.  The model scale must be part of a matching
analysis and carry uncertainty.

\subsection{Objective function}

At one truncation level, a provisional calibration objective is
\begin{equation}
 \chi_r^2(\theta_r)
 =\bigl(\bm y_{\cal}-\bm m_r(\theta_r)\bigr)^{\mathsf T}
 \bm C_{\cal}^{-1}
 \bigl(\bm y_{\cal}-\bm m_r(\theta_r)\bigr)
 +\cP_{\rm natural}(\theta_r)
 +\cP_{\rm renorm}(r),
 \label{eq:calibration-objective}
\end{equation}
with naturalness and renormalization-flow penalties grounded in the operator
power counting.  Volume~VI later replaces or embeds this objective in the full
shared posterior.

\begin{requirement}[Frozen prediction partition]
The calibration and withheld families are declared before the final parameter
optimization at each model version.  Promoting a failed holdout to calibration
creates a new version and requires a new independent holdout.
\label{req:frozen-holdout}
\end{requirement}

\section{Directed convergence program}
\label{sec:convergence}

\subsection{Multi-axis tower}

The convergence label is
\begin{equation}
 r=(K,N_{\max},b,\lambda,\epsilon_x,
 \cF,\text{color basis},\text{operator order},\text{numerics}).
 \label{eq:full-convergence-label}
\end{equation}
No single sequence explores every axis.  The program uses a designed set of
trajectories that separate:
\begin{itemize}
 \item longitudinal resolution;
 \item transverse UV and IR coverage;
 \item Hamiltonian resolution;
 \item Fock-sector content;
 \item induced-operator order;
 \item color/permutation completeness;
 \item quadrature and solver tolerances.
\end{itemize}

\subsection{Matched comparison}

If bases differ, a comparison map
\begin{equation}
 I_{r\to r'}:\Hil_r\to\Hil_{r'}
 \label{eq:comparison-map}
\end{equation}
and an operator matching map
\begin{equation}
 R_{r'\to r}:\cA_{r'}\to\cA_r
 \label{eq:operator-comparison}
\end{equation}
are supplied.  Observable convergence is assessed through
\begin{equation}
 \delta_i(r',r)
 =R_{r'\to r}\left[\cO_i^{(r')}\right]
 -\cO_i^{(r)}.
 \label{eq:convergence-residual}
\end{equation}

\subsection{Extrapolation model}

For a controlled trajectory, one may fit
\begin{equation}
 \cO_i^{(r)}
 =\cO_i^{(\infty)}
 +a_iK^{-p_i}
 +b_iN_{\max}^{-q_i}
 +c_i\exp(-d_iN_{\max})
 +e_iQ_{\cF}^{\nu_i}
 +\cdots,
 \label{eq:convergence-fit}
\end{equation}
but only terms justified by the regulator and basis analysis are admitted.
Alternative convergence forms form a model-discrepancy axis.

\subsection{Convergence categories}

The final report separates:
\begin{enumerate}
 \item eigenvalue and eigensolver error;
 \item basis extrapolation;
 \item Fock-sector truncation;
 \item regulator and renormalization trajectory;
 \item induced-operator truncation;
 \item rotational and Ward defects;
 \item current/operator matching;
 \item numerical transforms and integration.
\end{enumerate}
One category cannot be enlarged to hide failure in another.

\begin{proposition}[State-bundle convergence]
A state bundle converges at precision \(\eps\) only if the selected mass,
current, local-operator, sector-ledger, and nonzero-transfer overlap observables
all satisfy their category-specific tolerances along at least two independent
resolution trajectories.  Stability of the mass alone is insufficient.
\label{prop:bundle-convergence}
\end{proposition}

\section{The microscopic state bundle}
\label{sec:state-bundle}

\subsection{Purpose}

The eigensolver vector is not by itself a usable microscopic parent.  Every
accepted member is exported as a versioned \StateBundle containing the state,
Hamiltonian identity, operators, renormalization and convergence information,
and correlated proton/neutron structure.

\begin{definition}[Microscopic state bundle]
For member \(s\) and truncation level \(r\),
\begin{equation}
 \mathfrak B_{N,s}^{(r)}=
 \left(
 \begin{gathered}
 \mathfrak h_r,\mathfrak R_r,
 M_{N,s}^{2},|\Psi_{N,s}^{(r)}\rangle,
 \cP_{\nu}|\Psi_{N,s}^{(r)}\rangle,\\
 \text{phase convention},\text{isospin partner},
 \text{currents},\text{GTMD operators},\text{Wilson-ready vertices},\\
 \text{sector/OAM/color ledgers},
 \text{derivatives},\text{covariance},
 \text{convergence and defect reports}
 \end{gathered}
 \right).
 \label{eq:state-bundle}
\end{equation}
\end{definition}

\subsection{Required amplitude exports}

For each sector and nucleon helicity, the bundle exports sparse amplitudes
\begin{equation}
 \psi_{\nu,\Lambda}^{(r)}
 \left(
 \{x_i,\bm\kappa_{iT},\lambda_i,f_i,c_i\};\alpha
 \right)
 \label{eq:exported-amplitudes}
\end{equation}
with basis and continuum-evaluation adapters.  The multiplicity label \(\alpha\)
retains color, permutation, and recoupling identity.

The export must support:
\begin{itemize}
 \item evaluation at the initial and final recoil coordinates of Volume~II;
 \item active quark, antiquark, and gluon slot iteration;
 \item target and parton helicity matrices;
 \item sector-diagonal and declared induced/off-diagonal operator insertions;
 \item \(\Lz\), spin, and color interference ancestry;
 \item exact complex phase convention across proton/neutron and truncation
 members.
\end{itemize}

\subsection{Correlated proton and neutron}

The proton and neutron bundles share one microscopic member identifier and one
strong Hamiltonian parameter draw.  At the symmetric point,
\begin{equation}
 |n,s\rangle=\cU_I|p,s\rangle,
 \label{eq:isospin-bundle}
\end{equation}
with controlled CSB additions evaluated from the same state family.  Independent
proton and neutron parameter fits are not permitted in the central model.

\subsection{Ledgers}

Each bundle reports:
\begin{align}
 1&=\sum_{\nu}Z_{\nu,s}^{(r)},
 &Z_{\nu,s}^{(r)}&=\langle\Psi_{N,s}|P_\nu|\Psi_{N,s}\rangle,
 \label{eq:sector-ledger}\\
 1&=\sum_{a=q,\bar q,g}\avg{x}_{a,s}^{(r)}+\Delta_{P^+,s}^{(r)},
 \label{eq:momentum-ledger}\\
 1&=\sum_f\left(N_{f,s}-N_{\bar f,s}\right)/3,
 \label{eq:baryon-ledger}\\
 Q_N&=\sum_f e_f\left(N_{f,s}-N_{\bar f,s}\right)
 +\Delta_{Q,s}^{(r)},
 \label{eq:charge-ledger}\\
 \frac12&=\frac12\Delta\Sigma_s+L_{q,s}+\Delta G_s+L_{g,s}
 +\Delta_{J,s}^{(r)}.
 \label{eq:angular-ledger}
\end{align}
Sector probabilities are diagnostic coordinates; the total ledgers are
renormalization and acceptance quantities.

\subsection{Readiness statuses}

A bundle carries independent statuses:
\begin{longtable}{L{0.30\textwidth}L{0.60\textwidth}}
\toprule
Status & Meaning\\
\midrule
\endhead
\texttt{EIGENSTATE\_VALIDATED} & Hermiticity, residual, symmetry, color,
statistics, and normalization gates pass at one truncation.\\
{\footnotesize\ttfamily RENORMALIZATION\_\allowbreak TRAJECTORY\_\allowbreak VALIDATED} & Shared renormalization
conditions and selected observables stabilize over the required tower.\\
\texttt{CURRENT\_READY} & Charge, Ward, component, and induced-current gates
pass.\\
\texttt{GTMD\_OVERLAP\_READY} & Complete nonzero-transfer amplitude and
operator exports satisfy the Volume~II interface.\\
\texttt{WILSON\_READY} & Required one-gluon vertices, intermediate sectors,
phase conventions, and cut metadata satisfy Volume~III.\\
\texttt{NUCLEAR\_MATCHING\_READY} & Correlated proton/neutron members,
complete helicity matrices, current/operator matching, phase/soft metadata, and
covariance satisfy Volume~IV.\\
\texttt{LF\_TO\_QCD\_MATCHING\_READY} & Closed regulated operator basis and
matching metadata satisfy Volume~V.\\
\texttt{INFERENCE\_READY} & Differentiability, member identity, uncertainty
components, and frozen validation partitions satisfy Volume~VI.\\
\bottomrule
\end{longtable}
No status is inferred from another.

\begin{requirement}[Fail-closed handoff]
A downstream layer may consume only a state bundle carrying its explicit
readiness status.  Missing covariance, helicity blocks, phase data, matching
metadata, or convergence evidence cannot be replaced by defaults.
\label{req:fail-closed-handoff}
\end{requirement}

\section{Interface to common GTMD and Wilson calculations}
\label{sec:gtmd-interface}

\subsection{T-even overlap handoff}

For the regulated straight-link or zeroth-rescattering operator, the state
bundle supplies
\begin{equation}
 W_{a/N,\lambda'\lambda}^{[\Gamma],(r)}
 (x,\kT,\DT)
 =\sum_{\nu,j\in a}
 \int[\dd\Gamma_\nu]\,
 \psi_{\nu,\lambda'}^{(r)*}(\text{out})
 \cP_{j}^{[\Gamma],(r)}
 \psi_{\nu,\lambda}^{(r)}(\text{in}),
 \label{eq:state-to-gtmd}
\end{equation}
using the one authoritative recoil map and one state normalization.  Quark,
antiquark, and gluon correlators are different operators acting on the same
state, as required by the light-front overlap representation of form factors
and generalized parton correlators \cite{BrodskyDiehlHwang2001,Meissner2009,Lorce2011}.

\subsection{Wilson handoff}

The one-gluon Wilson pilot requires the same \(qqqg\) amplitudes and quark--gluon
vertices as the Hamiltonian.  The bundle therefore exports a vertex identity
\begin{equation}
 \mathfrak v_{qg,r}
 =\left(
 g_{qg,r},\text{color tensor},\text{helicity kernel},
 \text{regulator},\text{instantaneous partners},
 \text{counterterms}
 \right),
 \label{eq:vertex-identity}
\end{equation}
and the corresponding intermediate-state energies.  A Wilson kernel based on
a different coupling or gluon state is an external model, not a prediction of
the Hamiltonian.

\subsection{OAM and spin-orbit ancestry}

Each named interference result retains the amplitude blocks that generated it.
For a harmonic \(m\), the contribution has an ancestry record
\begin{equation}
 \mathfrak a_m=
 \left(
 \nu,\nu',\Lz,\Lz',\lambda_a,\lambda_a',
 \Lambda,\Lambda',m,\text{operator},\text{Wilson order}
 \right).
 \label{eq:oam-ancestry}
\end{equation}
Disabling either amplitude block must remove the corresponding harmonic.

\subsection{Operator matching remains later}

The state bundle is a regulated low-resolution object.  Even when its common
marginals close, the physical TMD remains
\begin{equation}
 \widetilde W_{a/N}^{\QCD}
 =\cZ_{\LF\to\QCD,a}
 \otimes_x
 \widetilde W_{a/N}^{\LF}
 +\Delta_{a}^{\mathrm{match}}.
 \label{eq:later-matching}
\end{equation}
No Volume~VII acceptance gate may be described as complete TMD matching or
evolution.

\section{Uncertainty and ensemble construction}
\label{sec:uncertainty}

\subsection{Microscopic member axes}

A state ensemble samples or scans:
\begin{equation}
 s=\left(
 \theta_H,\theta_{\ct},\theta_{\ind},
 r,\text{operator order},\text{solver},\text{numerics}
 \right)_s.
 \label{eq:member-id}
\end{equation}
All amplitudes, currents, GTMDs, Wilson kernels, and later nuclear quantities
derived from member \(s\) retain that identity.

\subsection{Separated uncertainty components}

Volume~VII reports:
\begin{enumerate}
 \item calibration/posterior uncertainty in shared parameters;
 \item basis and longitudinal-resolution extrapolation;
 \item Fock-sector truncation;
 \item induced-operator truncation and naturalness;
 \item regulator and renormalization trajectory;
 \item rotational, Ward, and PCAC defects;
 \item eigensolver and matrix-element numerics;
 \item state-tracking and near-degeneracy uncertainty;
 \item emulator uncertainty, if a surrogate is used.
\end{enumerate}
It does not yet include LF-to-QCD matching, TMD evolution, nuclear many-body,
or process-factorization uncertainty, except as future status fields.

\subsection{Common covariance}

For bundle observables \(\cO_i\),
\begin{equation}
 \Cov(\cO_i,\cO_j)
 =\mathbb E_s
 \left[
 (\cO_{i,s}-\bar\cO_i)
 (\cO_{j,s}-\bar\cO_j)
 \right].
 \label{eq:bundle-covariance}
\end{equation}
The covariance retains proton--neutron, quark--gluon, current--GTMD, and
cross-truncation correlations.  Fock-sector envelopes may not be combined as
independent errors if they are generated by the same parameter or omitted
operator.

\section{Software architecture}
\label{sec:software}

\subsection{Core objects}

\begin{longtable}{L{0.27\textwidth}L{0.65\textwidth}}
\toprule
Object & Responsibility\\
\midrule
\endhead
\texttt{HamiltonianResolution} & \(K,N_{\max},b,\lambda,\epsilon_x\), Fock
set, zero-mode policy, basis identity, ultraviolet/infrared interpretation.\\
\texttt{PartonBasisState} & Species, longitudinal and transverse modes,
helicity, flavor, color, and statistics identity.\\
\texttt{ColorSingletBasis} & Complete invariant tensors, multiplicities,
generator residuals, recoupling maps, and serialization.\\
\texttt{PermutationBasis} & Irreducible permutation labels and exact
identical-fermion antisymmetrizer.\\
\texttt{PhysicalFockBasis} & Tensor product followed by color, statistics,
quantum-number, and residual-gauge projection.\\
\Ham & Versioned block-sparse mass-squared operator with canonical, induced,
and counterterm terms.\\
\texttt{HamiltonianTerm} & Domain/codomain sectors, kernel, symmetry,
regulator, Hermitian partner, parameter ownership, and provenance.\\
\RenDatum & Bare/effective parameters, counterterms, induced coefficients,
conditions, schemes, operator maps, and remainder model.\\
\texttt{CurrentOperator} & One-body, instantaneous, exchange/induced, and
counterterm current blocks consistent with the Hamiltonian.\\
\texttt{PartonicOperatorBasis} & Decorated quark, antiquark, gluon, local,
and GTMD operator blocks.\\
\texttt{SparseEigensolver} & Matrix-free multiplication, preconditioning,
residuals, state/subspace tracking, and deterministic restart metadata.\\
\texttt{EigenstateMember} & Normalized state, phase convention, sector/OAM
ledgers, current fingerprints, and derivatives.\\
{\small\ttfamily Truncation\allowbreak Tower} & Directed levels, comparison/matching maps, refits, convergence
models, and model-discrepancy estimates.\\
\StateBundle & Complete downstream export and independent readiness statuses.\\
{\small\ttfamily Hamiltonian\allowbreak Validation\allowbreak Report} & Requirement coverage, benchmark
residuals, negative injections, source hashes, and acceptance decision.\\
\bottomrule
\end{longtable}

\subsection{Term interface}

Each Hamiltonian term implements at least:
\begin{verbatim}
source_sector
 target_sector
 symmetry_signature
 parameter_block
 regulator_identity
 apply(vector, resolution)
 adjoint_term_id
 matrix_element(bra, ket)
 derivative(parameter_id)
 provenance_node
 approximation_status
\end{verbatim}
A term without a source/target sector, Hermitian partner, or parameter owner
cannot enter the central Hamiltonian.

\subsection{Determinism and provenance}

Basis ordering, color and permutation recouplings, sparse matrix assembly,
random seeds, solver tolerances, and state phase conventions are serialized.
Generated matrices and state bundles carry content hashes.  A result is
reproducible only if its Hamiltonian and renormalization identities are
recoverable, not merely if a final CSV has the same values.

\subsection{Automatic differentiation}

Differentiable kernels are preferred, but sparse eigensolver derivatives may
use custom adjoints.  The implementation compares automatic/adjoint,
Hellmann--Feynman, and finite-difference derivatives on benchmark problems.
Non-differentiable basis or state-reordering operations are isolated from the
physical parameter map.

\section{Formal requirements}
\label{sec:requirements}

The stable requirement identifiers for this volume are:
\begin{longtable}{L{0.16\textwidth}L{0.75\textwidth}}
\toprule
ID & Requirement\\
\midrule
\endhead
\texttt{V7.HAM.001} & One versioned resolution-indexed microscopic Hamiltonian
generates the proton and neutron states.\\
\texttt{V7.HAM.002} & The Hamiltonian contains canonical kinetic, quark--gluon,
allowed gluon, and instantaneous terms with Hermitian partners.\\
\texttt{V7.HAM.003} & Induced confinement, chiral, and symmetry-restoration
operators carry resolution flow, provenance, and ablations.\\
\texttt{V7.FOCK.001} & The initial admissible common state contains explicit
\(qqq\), \(qqqg\), and light \(qqqq\bar q\) sectors.\\
\texttt{V7.FOCK.002} & Required \(\Lz=0,\pm1,\pm2\) blocks are present or
identified as controlled zero tests.\\
\texttt{V7.SYM.001} & Every sector uses a complete color-singlet basis with
multiplicity labels and generator tests.\\
\texttt{V7.SYM.002} & Full identical-fermion antisymmetry is imposed before
diagonalization.\\
\texttt{V7.SYM.003} & Exact kinematical quantum numbers and residual gauge
constraints are fail-closed basis types.\\
\texttt{V7.BASIS.001} & Longitudinal, transverse, center-of-mass, endpoint, and
zero-mode conventions are explicit.\\
\texttt{V7.REG.001} & Basis cutoffs, Hamiltonian resolution, endpoint regulator,
TMD rapidity scale, and numerical cutoffs are distinct types.\\
\texttt{V7.REN.001} & Renormalization conditions are reimposed at each tower
level and parameters carry ownership.\\
\texttt{V7.REN.002} & Fock-sector-dependent masses and vertices satisfy shared
pole and Ward conditions.\\
\texttt{V7.REN.003} & State, Hamiltonian, current, and GTMD operators transform
together under sector elimination or similarity change.\\
\texttt{V7.ISO.001} & Strong isospin is the reference symmetry; CSB and QED are
controlled operators.\\
\texttt{V7.CHI.001} & Sea flavor asymmetry arises from shared chiral/5q dynamics,
not independent antiquark inputs or large ad hoc mass splitting.\\
\texttt{V7.CUR.001} & Charge, Ward, axial/PCAC, EMT, and component-consistency
residuals are measured.\\
\texttt{V7.POIN.001} & Poincar\'e and rotational defects decrease under the
tower or are represented by controlled induced operators.\\
\texttt{V7.SOLVE.001} & Every eigenpair has a residual, state/subspace tracking,
and deterministic phase convention.\\
\texttt{V7.CONV.001} & Acceptance requires basis and Fock enlargement with
counterterm refitting and observable-level convergence.\\
\texttt{V7.CONV.002} & Mass convergence cannot substitute for current, GTMD,
Ward, or rotational convergence.\\
\texttt{V7.BUNDLE.001} & Proton and neutron state bundles retain correlated
member identity, amplitudes, operators, ledgers, covariance, and defects.\\
\texttt{V7.BUNDLE.002} & Downstream GTMD, Wilson, nuclear, matching, and
inference interfaces are independent fail-closed statuses.\\
\texttt{V7.CAL.001} & The calibration set is restricted and the withheld
families are frozen before final optimization.\\
\texttt{V7.PRED.001} & No named PDF/TMD/GTMD width, phase, or normalization is a
Hamiltonian parameter.\\
\texttt{V7.PROV.001} & Explicit sectors and their induced replacements are
mutually exclusive unless an overlap subtraction is present.\\
\texttt{V7.NUM.001} & Solver, quadrature, basis, and serialization residuals are
separate from physical truncation errors.\\
\texttt{V7.OUT.001} & The first central microscopic result is a versioned state
bundle, not a plot or independently fitted distribution.\\
\bottomrule
\end{longtable}

\section{Analytic and finite-dimensional benchmarks}
\label{sec:benchmarks}

\subsection{Benchmark H-A: free-sector spectrum}

For each sector, switch off all interactions and verify the free invariant-mass
spectrum, degeneracies, \(\Jz\), color, statistics, and center-of-mass
factorization.  The matrix-free and explicitly assembled matrices must agree on
small spaces.

\subsection{Benchmark H-B: one-vertex Hermiticity}

Retain one \(qqq\leftrightarrow qqqg\) vertex and its adjoint.  Verify
\begin{equation}
 \inner{\phi}{V_{3,4g}\psi}
 =\inner{V_{4g,3}\phi}{\psi}^{*}
 \label{eq:vertex-hermiticity-benchmark}
\end{equation}
for random symmetry-allowed states and reject a missing phase, color factor,
regulator, or adjoint block.

\subsection{Benchmark H-C: color and permutation completeness}

Construct all singlet multiplicities in \(qqq\), \(qqqg\), and
\(qqqq\bar q\).  Verify generator annihilation, orthonormality, recoupling
unitarity, and complete antisymmetry.  Deliberately omit one singlet or use the
wrong antiquark generator and require failure.

\subsection{Benchmark H-D: Fock-sector renormalization toy model}

Use an exactly solvable multi-sector Hamiltonian with a self-energy loop.  Fit
sector counterterms at two and three retained sectors.  Verify that the pole
mass and charge remain fixed, while bare/sector parameters flow.

\subsection{Benchmark H-E: Feshbach Hamiltonian and operator}

Extend the C4 finite model to a sector-coupled Hamiltonian, current, and norm
kernel.  Verify that
\begin{equation}
 \langle\Psi|\cO|\Psi\rangle
 =\langle\psi_P|N_{\eff}^{-1/2}\cO_{\eff}N_{\eff}^{-1/2}|\psi_P\rangle
 \label{eq:feshbach-benchmark}
\end{equation}
and that using only \(P\cO P\) fails.

\subsection{Benchmark H-F: Ward identity}

Use a finite light-front QED or Abelianized quark--gluon subproblem with the
same regulator and sector truncation to verify charge renormalization and the
regulated Ward identity.  The benchmark should reproduce the exact analytic
result where available and isolate the defect caused by omitting an
instantaneous term or vertex counterterm.

\subsection{Benchmark H-G: chiral pair sector}

In a reduced meson/nucleon model, verify that the chiral operator produces a
pseudoscalar bound-state channel, satisfies the declared PCAC relation, and
generates a positive-\(x\) antiquark distribution only when the five-parton
sector is active.  Switching off \(G_\chi\) or the sector must remove the effect.

\subsection{Benchmark H-H: rotational restoration}

Compare form factors or matrix elements extracted from independent current
components in increasing \(N_{\max}\) and sector spaces.  The angular or
component defect must decrease or be quantitatively attributed to an induced
operator with a convergent coefficient.

\subsection{Benchmark H-I: common-parent closure}

On the first nontrivial eigenstate, evaluate the regulated GTMD, TMD, GPD, PDF,
and local-current routes and verify the Volume~II commuting diagram.  No
TMD-specific width or normalization may be introduced.

\subsection{Benchmark H-J: state tracking}

Create an avoided crossing in a controlled block.  Verify that overlap and
current fingerprints track the physical state while an eigenvalue-ordering
rule fails.

\subsection{Benchmark H-K: confinement flow}

Compare induced-confinement and zero-confinement trajectories.  The benchmark
records masses, radii, sector probabilities, Ward defects, and coefficients as
resolution increases.  It must distinguish useful infrared acceleration from
a confining term that simply absorbs all omitted physics.

\subsection{Benchmark H-L: microscopic-member propagation}

Generate at least two correlated parameter members and verify that proton,
neutron, quark, antiquark, gluon, current, and GTMD outputs retain the same
member ID and covariance ancestry.

\section{Mandatory negative tests}
\label{sec:negative-tests}

The implementation must deliberately inject and reject at least the following
fault classes with stable identifiers and structured diagnostics:
\begin{enumerate}
 \item use \(P^+P^-\) instead of \(2P^+P^-\) under the declared convention;
 \item identify the oscillator scale \(b\) as a physical TMD width;
 \item reuse a TMD rapidity regulator as a Hamiltonian vertex cutoff;
 \item include \(V_{3,4g}\) without its Hermitian-conjugate block;
 \item change the vertex regulator only in one direction;
 \item omit a required instantaneous interaction;
 \item use a color-singlet \(qqq\) state times a free gluon;
 \item omit one independent color-singlet multiplicity;
 \item use the fundamental instead of anti-fundamental antiquark generator;
 \item violate identical-quark antisymmetry across a cluster partition;
 \item symmetrize an eigenvector only after diagonalization;
 \item count an explicit five-quark sector and its induced sea operator
 simultaneously;
 \item count explicit pionlike correlations and an unmatched hadronic pion
 addition simultaneously;
 \item use a large \(u-d\) effective-mass split to produce the sea asymmetry;
 \item use independent proton and neutron microscopic parameters;
 \item switch from effective to current masses in the highest sector without a
 renormalization condition;
 \item use one sector-independent mass counterterm when the benchmark requires
 FSDR;
 \item fit one independent vertex coefficient per named TMD;
 \item calibrate an induced confinement coefficient once and freeze it over the
 tower;
 \item declare convergence from a single \((K,N_{\max})\) point;
 \item declare convergence from the nucleon mass alone;
 \item compare bare wave-function coefficients from unrelated bases without a
 comparison map;
 \item identify a state only by eigenvalue order through an avoided crossing;
 \item omit center-of-mass factorization or the Lawson diagnostic;
 \item accept a state with an eigensolver residual above tolerance;
 \item use a current not transformed with the Hamiltonian;
 \item eliminate a sector in \(H\) but use only \(P\cO P\);
 \item satisfy charge normalization by a final scalar rescaling of all currents;
 \item suppress a Ward defect by fitting independent current components;
 \item use a chiral sea operator without a corresponding axial/PCAC test;
 \item use an endpoint regulator as a predicted small-\(x\) shape;
 \item omit the zero-mode policy from the Hamiltonian identity;
 \item label a sector probability as a regulator-independent observable;
 \item drop \(|\Lz|=2\) blocks and claim complete pretzelosity or tensor support;
 \item drop \(qqqg\) and claim a gluon or Wilson-line prediction;
 \item drop \(qqqq\bar q\) and populate antiquarks from a separate PDF;
 \item insert a phenomenological Sivers or Boer--Mulders curve into the state
 bundle;
 \item apply the C5 Wilson kernel with a coupling different from the Hamiltonian
 vertex;
 \item erase the color/permutation multiplicity label in an amplitude export;
 \item export incomplete target or parton helicity matrices;
 \item lose the complex phase convention between proton and neutron bundles;
 \item mix microscopic members when forming a ratio or asymmetry;
 \item mark a bundle \texttt{NUCLEAR\_MATCHING\_READY} without covariance;
 \item mark a bundle \texttt{LF\_TO\_QCD\_MATCHING\_READY} without a closed
 operator basis;
 \item mark a regulated overlap as a soft-subtracted TMD;
 \item use DGLAP evolution to define the Hamiltonian scale without matching;
 \item allow a pilot state or Hamiltonian to enter the accepted 216-route
 phenomenological production root;
 \item mutate one of the eight authoritative phenomenological artifacts;
 \item change a Volume~I--VI normative source during implementation;
 \item serialize two different Hamiltonian identities to the same version ID;
 \item use non-deterministic basis ordering without recording a permutation;
 \item report a covariance assembled from independently shuffled members;
 \item promote a failed holdout into calibration without a new model version;
 \item describe an induced operator as fundamental QCD without provenance;
 \item hide a large rotational or Ward defect inside a generic theory band;
 \item claim continuum QCD from a finite matrix with no regulator flow.
\end{enumerate}

\section{Staged research and implementation program}
\label{sec:stages}

The implementation stages below are formal work packages.  Their final Codex
identifiers should be assigned after C6 reports the actual active-gluon and
soft-overlap APIs.

\subsection{H0: basis, symmetry, and matrix-element spine}

\begin{enumerate}
 \item Implement \texttt{HamiltonianResolution}, one-particle labels, and
 intrinsic basis conventions.
 \item Generate complete color-singlet multiplicities and exact permutation
 bases for \(qqq\), \(qqqg\), and \(qqqq\bar q\).
 \item Implement center-of-mass factorization and exact quantum-number blocks.
 \item Implement canonical term interfaces and the free/one-vertex benchmarks.
 \item Keep every result validation-only and disconnected from production.
\end{enumerate}

Acceptance: Benchmarks H-A--H-C and all corresponding negative tests pass.
No physical calibration occurs.

\subsection{H1: valence renormalization benchmark}

\begin{enumerate}
 \item Construct the \(qqq\) Hamiltonian with canonical/effective interactions.
 \item Implement induced-confinement and zero-confinement trajectories.
 \item Implement electromagnetic current and center-of-mass/current tests.
 \item Solve several \((K,N_{\max},b,\lambda)\) points and refit the shared
 conditions.
 \item Export a valence-only benchmark bundle explicitly marked incomplete.
\end{enumerate}

Acceptance: mass, charge, selected form factor, current-component, and
rotational diagnostics exhibit controlled flow.  No sea/gluon claim is made.

\subsection{H2: dynamical-gluon state}

\begin{enumerate}
 \item Add \(qqqg\) color/permutation sectors and canonical
 \(qqq\leftrightarrow qqqg\) vertices.
 \item Implement Fock-sector-dependent mass and vertex renormalization.
 \item Add instantaneous partners and the Abelianized Ward benchmark.
 \item Export gluon helicity, OAM, and field-strength operator blocks.
 \item Reconnect the C5/C6 Wilson pilot only through the Hamiltonian vertex and
 state bundle.
\end{enumerate}

Acceptance: H-D, H-F, H-H, and the gluon portions of H-I pass over a basis and
sector tower.

\subsection{H3: light sea and chiral dynamics}

\begin{enumerate}
 \item Add fully antisymmetrized \(u\bar u\) and \(d\bar d\) five-parton
 sectors.
 \item Implement the shared chiral operator basis and axial/PCAC current.
 \item Calibrate the minimal chiral conditions and withhold the broader sea
 shape/asymmetry.
 \item Implement explicit/induced sea exclusion and overlap subtraction.
 \item Export positive-\(x\) antiquark helicity matrices from the state.
\end{enumerate}

Acceptance: H-C, H-E, H-G, H-I, and sea-flavor withheld tests pass with a
controlled truncation estimate.

\subsection{H4: common microscopic GTMD bundle}

\begin{enumerate}
 \item Evaluate nonzero-\(\DT\) quark, antiquark, and gluon overlaps from the
 same state.
 \item Verify TMD/GPD/PDF/current and OAM reductions.
 \item Export complete proton/neutron helicity matrices and microscopic member
 covariance.
 \item Produce the first \texttt{GTMD\_OVERLAP\_READY} state bundle.
 \item Compare with, but do not absorb, the accepted phenomenological boundary.
\end{enumerate}

\subsection{H5: Wilson, nuclear, matching, and inference readiness}

\begin{enumerate}
 \item Demonstrate that the Wilson kernel uses the Hamiltonian gluon field,
 vertex, and intermediate states.
 \item Complete the current/operator basis needed by LF-to-QCD matching.
 \item Generate correlated proton/neutron members and covariance for nuclear
 composition.
 \item Provide derivatives and exact-solver benchmarks for Volume~VI inference.
 \item Freeze an independent family of nucleon observables as withheld
 prediction gates.
\end{enumerate}

This stage does not itself perform the full deuteron calculation, QCD evolution,
or global fit; it certifies the microscopic input.

\section{Risk register and safeguards}
\label{sec:risks}

\begin{longtable}{L{0.23\textwidth}L{0.34\textwidth}L{0.35\textwidth}}
\toprule
Risk & Failure mode & Required safeguard\\
\midrule
\endhead
Finite-matrix realism & A smooth spectrum and plausible PDFs are mistaken for
renormalized QCD & Publish regulator, counterterm, basis, Fock, Ward, rotational,
and withheld-observable convergence together\\
Induced-potential overreach & Confinement or chiral operators absorb all missing
physics and become one flexible fit per channel & Shared operator basis,
naturalness, resolution flow, zero-operator ablations, and withheld channels\\
Sector-dependent patchwork & Different masses or couplings are assigned to each
sector without renormalization conditions & FSDR datum, Ward constraints,
sector-promotion tests, and parameter-ownership audit\\
Incomplete antisymmetry & Ordered cluster bases double count or misnormalize
five-parton states & Exact permutation projection before diagonalization and
multiplicity tests\\
Color incompleteness & A single convenient singlet misses independent color
channels & Generate the full invariant subspace and test recoupling completeness\\
Zero-mode blindness & Omitting \(k^+=0\) is treated as harmless & Typed zero-mode
policy, induced remainder, endpoint studies, and later zero-mode-resolved tower\\
Confinement double counting & Explicit gluon/chiral sectors duplicate induced
operators & Provenance exclusions, overlap subtraction, and Feshbach comparison\\
Current inconsistency & State and current are taken from different Hamiltonians &
Gauge/transform Hamiltonian and operators together; enforce Ward and component
tests\\
Rotational overfitting & Independent current coefficients force angular
conditions at one resolution & Enlarge basis/Fock space first; restrict induced
operators and test withheld components\\
False convergence & One basis path or mass plateau hides large current/GTMD
defects & Multi-axis trajectories and state-bundle convergence criterion\\
State misidentification & Avoided crossings swap the nucleon eigenvector &
Subspace tracking with overlap, currents, and quantum-number fingerprints\\
Sea by mass splitting & Large \(u-d\) masses imitate flavor asymmetry & Strong
isospin reference plus chiral five-parton dynamics and controlled CSB\\
Scale conflation & Hamiltonian, basis, TMD rapidity, and evolution scales are
interchanged & Strong typed scale identities and negative injections\\
Premature production & Validation state silently replaces the accepted boundary &
Disjoint provenance root and explicit readiness statuses\\
Computational explosion & Full sectors become intractable before convergence &
Symmetry blocks, matrix-free solvers, distributed contractions, reduced-basis
emulators, and staged promotion\\
\bottomrule
\end{longtable}

\section{Final acceptance criteria}
\label{sec:acceptance}

The first microscopic nucleon state may be used as the central common parent
only when all of the following hold:
\begin{enumerate}
 \item One versioned RSA-BLFQ Hamiltonian and renormalization trajectory
 generate correlated proton and neutron states.
 \item The accepted state contains explicit \(qqq\), \(qqqg\), and light
 \(qqqq\bar q\) sectors.
 \item Every sector is color singlet, fully antisymmetrized where required, and
 free of unacceptable center-of-mass contamination.
 \item The Hamiltonian is Hermitian and every eigenpair satisfies the declared
 residual tolerance.
 \item The regulator, zero-mode, instantaneous-kernel, and endpoint policies are
 explicit and reproducible.
 \item Sector-dependent masses, vertices, and counterterms satisfy shared
 renormalization conditions and Ward constraints.
 \item Induced confinement and chiral operators have explicit provenance,
 natural coefficients, resolution flow, and successful ablations.
 \item Strong isospin and controlled CSB produce correlated proton/neutron
 flavor structure.
 \item The state contains the \(\Lz\) and helicity blocks required by the
 declared GTMD and Wilson projections.
 \item Charge, baryon number, plus momentum, and angular-momentum ledgers close
 within separate tolerances.
 \item Electromagnetic, axial/PCAC, and EMT operators are consistent with the
 Hamiltonian and sector elimination.
 \item Poincar\'e, rotational, and current-component defects decrease under the
 truncation tower or are accounted for by controlled induced operators.
 \item Masses and selected calibration observables are stable after both basis
 and Fock-sector enlargement with counterterm refitting.
 \item At least one electromagnetic, one spin/chiral, one partonic, and one
 nonzero-transfer observable family excluded from calibration is successfully
 predicted.
 \item Quark, antiquark, and gluon GTMDs arise from the same state and their
 regulated TMD, GPD, PDF, and current reductions commute.
 \item The Wilson kernel uses the same gluon field, coupling, color structures,
 and intermediate-state spectrum as the Hamiltonian.
 \item Removing every presently phenomenological named-TMD input leaves the
 microscopic calculation defined.
 \item The complete \StateBundle retains amplitudes, operators, phases,
 proton/neutron member identity, covariance, derivatives, convergence, and
 defect manifests.
 \item Downstream readiness statuses are assigned independently and only after
 their corresponding gates pass.
 \item All analytic benchmarks, property tests, negative injections,
 provenance exclusions, and numerical regression gates pass.
\end{enumerate}

\begin{decision}[Meaning of the first accepted result]
Passing \cref{sec:acceptance} establishes a \emph{microscopic nucleon state
bundle admissible for common GTMD prediction}.  It does not yet establish a
physical soft-subtracted TMD, a complete deuteron state, a process cross
section, or a global posterior.  Those claims require the independent matching,
nuclear, process, and inference gates of Volumes~IV--VI.
\label{dec:meaning-accepted-result}
\end{decision}

\section{Recommended implementation boundary after C6}
\label{sec:next-package}

C6 is expected to complete the validation-only active-gluon ordered-link
\(f/d\) rescattering and analytic soft-overlap accounting.  The next software
package should then begin Volume~VII at H0 rather than attempting the full
Hamiltonian at once.

The recommended boundary is:
\begin{quote}
\textbf{C7/H0: symmetry-complete microscopic basis and Hamiltonian-term spine.}
Implement the resolution identity, complete color-singlet and permutation
bases for \(qqq\), \(qqqg\), and \(qqqq\bar q\), center-of-mass factorization,
free invariant-mass operator, typed canonical term interfaces, one-vertex
Hermiticity, and deterministic serialization.  Use analytic matrices only.
Do not calibrate a nucleon, modify production artifacts, or claim a microscopic
eigenstate until these basis and algebraic gates pass.
\end{quote}

The exact Codex prompt must be written after the C6 report so it can consume the
actual ordered-link, color-channel, soft-overlap, provenance, and normative
source APIs.

\section{Conclusion}
\label{sec:conclusion}

The architecture developed in Volumes~I--VI becomes predictive only after the
central amplitudes cease to be supplied by unrelated PDF, TMD, phase, and
nuclear models.  Volume~VII makes the primary dynamical choice needed for that
transition.  RSA-BLFQ is selected because it can represent explicit quarks,
antiquarks, gluons, color and permutation multiplicities, helicity and orbital
blocks, sector-changing QCD vertices, and nonzero-transfer overlaps in one
sparse Hamiltonian framework.  Its use is deliberately stricter than adopting
an available finite BLFQ parameterization.  The Hamiltonian is a
resolution-indexed EFT whose canonical terms, induced confinement and chiral
operators, sector-dependent counterterms, currents, and partonic operators
must flow together and converge at the observable level.

The first scientifically meaningful output is not a new atlas of smooth TMD
curves.  It is a correlated proton/neutron state bundle carrying the complete
microscopic amplitudes, operator identities, phases, ledgers, covariance, and
convergence record required by the GTMD, Wilson, nuclear, matching, and
inference layers.  Only after that bundle succeeds on withheld observables may
the project replace the accepted phenomenological boundary with a genuinely
generative microscopic parent.

\appendix

\section{Detailed sector table}
\label{app:sectors}

\begingroup
\setlength{\tabcolsep}{3pt}
\begin{longtable}{L{0.11\textwidth}L{0.16\textwidth}L{0.19\textwidth}L{0.20\textwidth}L{0.18\textwidth}}
\toprule
Sector & Minimum color structure & Required orbital/helicity support & Primary physics & First acceptance evidence\\
\midrule
\endhead
\(qqq\) & Unique baryon singlet, exact permutation coupling &
\(\Lz=0,\pm1,\pm2\) as allowed by \(N_{\max}\); nucleon
\(\Lambda=\pm\frac12\) & Valence mass, charge, helicity, transversity,
baseline OAM & Spectrum, charge, form factors, current components,
color/statistics\\
\(qqqg\) & All singlets from three-quark octet multiplicities coupled to
adjoint & Gluon \(\lambda_g=\pm1\), quark helicities, required \(\Lz\)
interferences & Gluon PDF/helicity/OAM, self-energy, Wilson cut, genuine
quark--gluon terms & Ward toy model, gluon momentum, field-strength projectors,
Wilson consistency\\
\(qqqu\bar u\) & All five-body singlet multiplicities, anti-fundamental
antiquark & Active \(u\), active \(\bar u\), cluster and non-cluster orbital
blocks & Light sea, chiral correlations, antiquark GTMDs & Antisymmetry, number,
charge, momentum, PCAC, sea holdout\\
\(qqqd\bar d\) & Same as above & Same, with correlated isospin coupling &
\(\bar d-\bar u\), pionlike correlations & Isospin, chiral conditions, sea
flavor holdout\\
\(qqqgg\) & Complete singlets including multiple adjoint recouplings &
Two-gluon helicity/color and higher \(\Lz\) support & Non-Abelian Wilson order,
gluon self-interactions, Ward closure & Promoted when H2/H3 defects demand it\\
\(qqqq\bar qg\) & Complete six-parton singlets and antisymmetry & Combined sea,
gluon, and OAM blocks & Mixed sea--gluon evolution and higher Wilson states &
Promoted by operator closure or truncation remainder\\
\bottomrule
\end{longtable}
\endgroup

\section{Hamiltonian identity schema}
\label{app:ham-schema}

\begin{verbatim}
{
  "hamiltonian_id": "rsa_blfq_nucleon_v...",
  "light_front_convention": "vpm=(v0+-v3)/sqrt(2)",
  "gauge": "Aplus=0",
  "resolution": {
    "K": "...",
    "Nmax": "...",
    "b_GeV": "...",
    "lambda_vertex_GeV": "...",
    "endpoint_regulator": "...",
    "zero_mode_policy_id": "..."
  },
  "fock_sectors": ["qqq", "qqqg", "qqquubar", "qqqddbar"],
  "color_basis_id": "...",
  "permutation_basis_id": "...",
  "canonical_terms": ["kinetic", "qg", "instantaneous", "..."],
  "induced_terms": ["confinement", "chiral", "rotation_restore"],
  "counterterm_basis": ["..."],
  "renormalization_conditions": ["..."],
  "current_basis_id": "...",
  "partonic_operator_basis_id": "...",
  "provenance_graph_id": "...",
  "source_hashes": {"...": "sha256"},
  "version": "..."
}
\end{verbatim}

\section{State-bundle schema}
\label{app:bundle-schema}

\begin{verbatim}
{
  "bundle_id": "...",
  "microscopic_member_id": "...",
  "nucleon": "proton|neutron",
  "isospin_partner_bundle_id": "...",
  "hamiltonian_id": "...",
  "renormalization_datum_id": "...",
  "mass_GeV": "...",
  "eigensolver_residual": "...",
  "phase_convention": "...",
  "amplitudes": {
    "qqq": "sparse reference",
    "qqqg": "sparse reference",
    "qqquubar": "sparse reference",
    "qqqddbar": "sparse reference"
  },
  "sector_ledger": {"...": "..."},
  "momentum_ledger": {"...": "..."},
  "spin_oam_ledger": {"...": "..."},
  "current_operator_ids": ["..."],
  "gtmd_operator_basis_id": "...",
  "wilson_vertex_id": "...",
  "derivative_manifest_id": "...",
  "covariance_member_id": "...",
  "convergence_report_id": "...",
  "defect_report_id": "...",
  "readiness": {
    "eigenstate_validated": true,
    "renormalization_trajectory_validated": false,
    "current_ready": false,
    "gtmd_overlap_ready": false,
    "wilson_ready": false,
    "nuclear_matching_ready": false,
    "lf_to_qcd_matching_ready": false,
    "inference_ready": false
  }
}
\end{verbatim}

\section{Renormalization-trajectory manifest}
\label{app:trajectory-schema}

\begin{verbatim}
{
  "trajectory_id": "...",
  "fixed_physical_conditions": ["M_N", "F1p(0)", "..."],
  "levels": [
    {
      "resolution_id": "...",
      "fock_set": ["..."],
      "fitted_parameters": {"...": "..."},
      "counterterms": {"...": "..."},
      "induced_coefficients": {"...": "..."},
      "calibration_residuals": {"...": "..."},
      "withheld_observables": {"...": "..."},
      "ward_defects": {"...": "..."},
      "rotational_defects": {"...": "..."},
      "operator_matching_residuals": {"...": "..."}
    }
  ],
  "comparison_maps": ["..."],
  "convergence_models": ["..."],
  "accepted_precision": {"...": "..."}
}
\end{verbatim}

\section{Compact adoption checklist}
\label{app:checklist}

Before any Hamiltonian result is described as microscopic, answer:
\begin{enumerate}
 \item Is the complete Hamiltonian identity recoverable?
 \item Are all retained sectors color complete and fully antisymmetrized?
 \item Are the zero-mode and endpoint policies explicit?
 \item Are canonical vertices accompanied by instantaneous and Hermitian
 partners?
 \item Do induced interactions carry provenance, flow, and ablations?
 \item Are counterterms tied to shared renormalization conditions?
 \item Are Hamiltonian and operators transformed together?
 \item Do Ward, PCAC, EMT, and rotational residuals decrease?
 \item Is the nucleon tracked by state/subspace identity across the tower?
 \item Has at least one basis and one Fock enlargement been completed?
 \item Do common GTMD reductions close without named-function repairs?
 \item Are proton and neutron correlated members of one strong Hamiltonian?
 \item Is every downstream readiness status independently justified?
 \item Are calibration and holdout families frozen?
 \item Can the calculation run with all phenomenological named-TMD inputs
 removed?
\end{enumerate}

\small
\begin{thebibliography}{99}

\bibitem{Dirac1949}
P.~A.~M.~Dirac,
``Forms of relativistic dynamics,''
Rev. Mod. Phys. \textbf{21}, 392 (1949).

\bibitem{BrodskyPauliPinsky1998}
S.~J.~Brodsky, H.-C.~Pauli, and S.~S.~Pinsky,
``Quantum chromodynamics and other field theories on the light cone,''
Phys. Rept. \textbf{301}, 299 (1998),
arXiv:hep-ph/9705477.

\bibitem{Vary2010}
J.~P.~Vary et al.,
``Hamiltonian light-front field theory in a basis function approach,''
Phys. Rev. C \textbf{81}, 035205 (2010),
arXiv:0905.1411.

\bibitem{Zhao2014}
X.~Zhao, H.~Honkanen, P.~Maris, J.~P.~Vary, and S.~J.~Brodsky,
``Electron g-2 in light-front quantization,''
Phys. Lett. B \textbf{737}, 65 (2014),
arXiv:1402.4195.

\bibitem{Xu2021}
S.~Xu et al. (BLFQ Collaboration),
``Nucleon structure from basis light-front quantization,''
Phys. Rev. D \textbf{104}, 094036 (2021),
arXiv:2108.03909.

\bibitem{Xu2023}
S.~Xu et al. (BLFQ Collaboration),
``Quark and gluon spin and orbital angular momentum in the proton,''
Phys. Rev. D \textbf{108}, 094002 (2023),
arXiv:2209.08584.

\bibitem{Xu2025}
S.~Xu, Y.~Liu, C.~Mondal, J.~Lan, X.~Zhao, Y.~Li, and J.~P.~Vary,
``Towards a first principles light-front Hamiltonian for the nucleon,''
Phys. Lett. B \textbf{867}, 139599 (2025),
arXiv:2408.11298.

\bibitem{Mondal2025}
C.~Mondal, S.~Xu, Y.~Liu, J.~Lan, Y.~Li, X.~Zhao, and J.~P.~Vary,
``Basis light-front quantization: Advancing a first principles approach for the
nucleon,''
arXiv:2503.21372 (2025).

\bibitem{Karmanov2005}
J.-F.~Mathiot, V.~A.~Karmanov, and A.~V.~Smirnov,
``Non-perturbative renormalization in light-front dynamics with Fock space
truncation,''
Nucl. Phys. B Proc. Suppl. \textbf{161}, 160 (2006),
arXiv:hep-th/0510230.

\bibitem{Karmanov2008}
V.~A.~Karmanov, J.-F.~Mathiot, and A.~V.~Smirnov,
``Systematic renormalization scheme in light-front dynamics with Fock space
truncation,''
Phys. Rev. D \textbf{77}, 085028 (2008),
arXiv:0801.4507.

\bibitem{Karmanov2012}
V.~A.~Karmanov, J.-F.~Mathiot, and A.~V.~Smirnov,
``Nonperturbative calculation of the anomalous magnetic moment in the Yukawa
model within truncated Fock space,''
Phys. Rev. D \textbf{86}, 085006 (2012),
arXiv:1204.3257.

\bibitem{Wilson1994}
K.~G.~Wilson, T.~S.~Walhout, A.~Harindranath, W.-M.~Zhang,
R.~J.~Perry, and S.~D.~Glazek,
``Nonperturbative QCD: A weak-coupling treatment on the light front,''
Phys. Rev. D \textbf{49}, 6720 (1994),
arXiv:hep-th/9401153.

\bibitem{GlazekWilson1993}
S.~D.~Glazek and K.~G.~Wilson,
``Renormalization of Hamiltonians,''
Phys. Rev. D \textbf{48}, 5863 (1993).

\bibitem{Wegner1994}
F.~Wegner,
``Flow-equations for Hamiltonians,''
Ann. Phys. (Leipzig) \textbf{3}, 77 (1994).

\bibitem{GlazekPerry1992}
S.~D.~Glazek and R.~J.~Perry,
``Special example of relativistic Hamiltonian field theory,''
Phys. Rev. D \textbf{45}, 3740 (1992).

\bibitem{Burkardt1998}
M.~Burkardt,
``Finiteness conditions for light-front Hamiltonians,''
Phys. Rev. D \textbf{58}, 096015 (1998),
arXiv:hep-th/9805088.

\bibitem{Feshbach1958}
H.~Feshbach,
``Unified theory of nuclear reactions,''
Ann. Phys. \textbf{5}, 357 (1958).

\bibitem{Okubo1954}
S.~Okubo,
``Diagonalization of Hamiltonian and Tamm--Dancoff equation,''
Prog. Theor. Phys. \textbf{12}, 603 (1954).

\bibitem{LeeSuzuki1980}
S.~Y.~Lee and K.~Suzuki,
``The effective interaction of two nucleons in the \(sd\) shell,''
Phys. Lett. B \textbf{91}, 173 (1980).

\bibitem{Li2016}
Y.~Li, P.~Maris, X.~Zhao, and J.~P.~Vary,
``Heavy quarkonium in a holographic basis,''
Phys. Lett. B \textbf{758}, 118 (2016),
arXiv:1509.07212.

\bibitem{Cao2026}
X.~Cao, S.~Cheng, Y.~Duan, Y.~Li, S.~Xu, and X.~Zhao,
``Non-perturbative flavor asymmetry in the nucleon and deuteron: The
light-front Hamiltonian effective field theory approach,''
J. Subatomic Part. Cosmol. \textbf{5}, 100329 (2026),
arXiv:2601.13567.

\bibitem{Wu2026}
X.~Wu et al.,
``Nuclear matter and proton parton distributions in a light-front Hamiltonian
framework,''
arXiv:2605.31030 (2026).

\bibitem{BrodskyDiehlHwang2001}
S.~J.~Brodsky, M.~Diehl, and D.~S.~Hwang,
``Light-cone wavefunction representation of deeply virtual Compton scattering,''
Nucl. Phys. B \textbf{596}, 99 (2001),
arXiv:hep-ph/0009254.

\bibitem{Meissner2009}
S.~Mei{\ss}ner, A.~Metz, and M.~Schlegel,
``Generalized parton correlation functions for a spin-1/2 hadron,''
JHEP \textbf{08}, 056 (2009),
arXiv:0906.5323.

\bibitem{Lorce2011}
C.~Lorc\'e and B.~Pasquini,
``Quark Wigner distributions and orbital angular momentum,''
Phys. Rev. D \textbf{84}, 014015 (2011),
arXiv:1106.0139.

\bibitem{Collins2011}
J.~Collins,
\emph{Foundations of Perturbative QCD},
Cambridge University Press (2011).

\end{thebibliography}

\end{document}
